\documentclass[11pt]{article}
\usepackage[a4paper,margin=22mm]{geometry}
\usepackage[no-math]{fontspec}
\setmainfont{lmroman10-regular.otf}[BoldFont=lmroman10-bold.otf,ItalicFont=lmroman10-italic.otf,BoldItalicFont=lmroman10-bolditalic.otf]
\usepackage{amsmath,amssymb,amsthm,mathtools,graphicx,xcolor,hyperref,xurl,xspace}
\setlength{\emergencystretch}{4em}
\SetMathAlphabet{\mathbf}{normal}{OT1}{cmr}{bx}{n}
\SetMathAlphabet{\mathbf}{bold}{OT1}{cmr}{bx}{n}
\newtheorem{theo}{Theorem}[section]
\newtheorem{lemm}[theo]{Lemma}
\newtheorem{lem}[theo]{Lemma}
\newtheorem{lemma}[theo]{Lemma}
\newtheorem{prop}[theo]{Proposition}
\newtheorem{coro}[theo]{Corollary}
\newtheorem{corol}[theo]{Corollary}
\newtheorem{conj}[theo]{Conjecture}
\newtheorem{Proposition}[theo]{Proposition}
\newtheorem{theoreme}[theo]{Theorem}
\theoremstyle{definition}
\newtheorem{defi}[theo]{Definition}
\newtheorem{exam}[theo]{Example}
\newtheorem{exem}[theo]{Example}
\newtheorem{rema}[theo]{Remark}
\newtheorem{nota}[theo]{Notation}
\newtheorem*{rema*}{Remark}
\newtheorem*{defi*}{Definition}
\newtheorem*{theo*}{Theorem}
\newtheorem*{lemm*}{Lemma}
\newtheorem*{prop*}{Proposition}
\newcommand{\enoncename}{}
\newtheorem{corpusenonce}[theo]{\enoncename}
\newenvironment{enonce}[1]{\renewcommand{\enoncename}{#1}\begin{corpusenonce}}{\end{corpusenonce}}
\newenvironment{enonce*}[1]{\par\medskip\noindent\textbf{#1.}\itshape}{\par\medskip}
\newenvironment{proofc}{\begin{proof}}{\end{proof}}
\newcommand{\Mc}{\mathcal{M}}
\newcommand{\pointir}{.\ ---\ }
\providecommand{\mathof}[1]{\mathrm{#1}}

\numberwithin{equation}{section}
\numberwithin{figure}{section}
\numberwithin{table}{section}
\def\endappendix{}
%~Mouliné par MaN_auto v.0.23.0 2020-08-12 15:12:34


\usepackage{tikz}
\usepackage{tikz-cd}
%\usepackage{enumitem}
%\usepackage{etoolbox}
%\usepackage[explicit]{titlesec}
%\usepackage{xspace}



\AtBeginDocument{
\newtheorem{constr}[theo]{Construction}
\newtheorem{thmX}{Theorem}
\renewcommand{\thethmX}{\Alph{thmX}} }

\DeclareMathOperator*{\fibprod}{\times}
\DeclarePairedDelimiter\abs{\lvert}{\rvert}
%\DeclarePairedDelimiter\norm{\lVert}{\rVert}



\usetikzlibrary{matrix}

\tikzset{commutative diagrams/.cd, arrow style=tikz, diagrams={>=latex}}




%%ENSUREMATH
\newcommand{\sC}{\ensuremath{\mathcal{C}}\xspace}
\newcommand{\sE}{\ensuremath{\mathcal{E}}\xspace}
\newcommand{\sF}{\ensuremath{\mathcal{F}}\xspace}
\newcommand{\sG}{\ensuremath{\mathcal{G}}\xspace}
\newcommand{\sI}{\ensuremath{\mathcal{I}}\xspace}
\newcommand{\sL}{\ensuremath{\mathcal{L}}\xspace}
\newcommand{\sN}{\ensuremath{\mathcal{N}}\xspace}
\newcommand{\sO}{\ensuremath{\mathcal{O}}\xspace}
\newcommand{\bA}{\ensuremath{\mathbf{A}}\xspace}
\newcommand{\bB}{\ensuremath{\mathbf{B}}\xspace}
\newcommand{\bC}{\ensuremath{\mathbf{C}}\xspace}
\newcommand{\bD}{\ensuremath{\mathbf{D}}\xspace}
\newcommand{\bG}{\ensuremath{\mathbf{G}}\xspace}
\newcommand{\bP}{\ensuremath{\mathbf{P}}\xspace}
\newcommand{\bZ}{\ensuremath{\mathbf{Z}}\xspace}
\newcommand{\bV}{\ensuremath{\mathbf{V}}\xspace}



\newcommand{\colim}{\varinjlim}
\renewcommand{\lim}{\varprojlim}

\newcommand{\initialobj}{\varnothing}
\renewcommand{\setminus}{\smallsetminus}

\newcommand{\hook}{\hookrightarrow}
\newcommand{\isoto}{\xrightarrow{\sim}}
\renewcommand{\ge}{\geqslant}
\renewcommand{\le}{\leqslant}




\newcommand{\Spc}{\mrm{Spc}}
\newcommand{\SCRing}{\mrm{SCRing}}
\newcommand{\Mod}{\mrm{Mod}}
\newcommand{\perf}{{\mrm{perf}}}


\newcommand{\A}{\bA}
\newcommand{\red}{{\mrm{red}}}
\newcommand{\bDelta}{\mathbf{\Delta}}

\newcommand{\cl}{{\mrm{cl}}}
\newcommand{\idem}{{\mrm{idem}}}

\renewcommand{\P}{\bP}

%%MATHOP%%%%%
\DeclareMathOperator{\Maps}{Maps}
\DeclareMathOperator{\Sym}{Sym}
\DeclareMathOperator{\rk}{rk}
\DeclareMathOperator{\Spec}{Spec}
\DeclareMathOperator{\cn}{cn}
\DeclareMathOperator{\Qcoh}{Qcoh}
\DeclareMathOperator{\Perf}{Perf}
\DeclareMathOperator{\K}{K}

\DeclareMathOperator{\Bl}{Bl}
\DeclareMathOperator{\KH}{KH}
\DeclareMathOperator{\id}{id}
\DeclareMathOperator{\op}{op}

\newcommand{\term}[1]{#1\xspace}
%\newcommand{\scr}{\term{simplicial commutative ring}}
%\newcommand{\scrs}{\term{simplicial commutative rings}}

%\newcommand{\inftyCat}{\term{$\infty$-category}}
%\newcommand{\inftyCats}{\term{$\infty$-categories}}

%\newcommand{\das}{\term{derived algebraic space}}
%\newcommand{\dass}{\term{derived algebraic spaces}}

%\newcommand{\as}{\term{algebraic space }}
%\newcommand{\ass}{\term{algebraic spaces }}


\newcommand{\mrm}[1]{\ensuremath{\mathrm{#1}}\xspace}








%%%%%%%%%%%%%%%%%%%%%%%%%%%%%%%%%%%%%%%%%%%%%%%%%
\graphicspath{{./figures/}}

\newcommand*{\mk}{\mkern -1mu}
\newcommand*{\Mk}{\mkern -2mu}
\newcommand*{\mK}{\mkern 1mu}
\newcommand*{\MK}{\mkern 2mu}

\hypersetup{urlcolor=purple, linkcolor=blue, citecolor=red}

\newcommand*{\romanenumi}{\renewcommand*{\theenumi}{\roman{enumi}}}
\newcommand*{\Romanenumi}{\renewcommand*{\theenumi}{\Roman{enumi}}}
\newcommand*{\alphenumi}{\renewcommand*{\theenumi}{\alph{enumi}}}
\newcommand*{\Alphenumi}{\renewcommand*{\theenumi}{\Alph{enumi}}}
\let\oldtilde\tilde
\renewcommand*{\tilde}[1]{\mathchoice{\widetilde{#1}}{\widetilde{#1}}{\oldtilde{#1}}{\oldtilde{#1}}}
%%%%%%%%%%%%%%%%%%%%%%%%%%%%%%%%%%%%%%%%%%%%%%%%%
\begin{document}
\section*{\raggedright 2634: Semi-orthogonal decomposition of perfect complexes on blow-ups of derived Artin stacks in quasi-smooth closed centres}
\noindent\textbf{Author:} Adeel A. Khan\par\noindent\textbf{Source:} Algebraic K-theory of quasi-smooth blow-ups and cdh descent. Annales Henri Lebesgue 3 (2020), publisher version, DOI 10.5802/ahl.55. \url{https://ahl.centre-mersenne.org/item/AHL_2020__3__1091_0/}
\par\noindent\textbf{License:} CC BY 4.0, \url{https://creativecommons.org/licenses/by/4.0/}. The original publisher PDF has the explicit license notice. See the saved source/license records.
\par\noindent\textbf{Editorial material note (not author text):} The selected problem is Theorem C, the ordered semi-orthogonal decomposition of perfect complexes on a derived Artin-stack blow-up in a quasi-smooth closed centre of virtual codimension at least one. The complete author article body retains all local derived projective-bundle, virtual-Cartier splitting/duality, blow-up full-faithfulness, orthogonality, generation and perfect-preservation arguments. Existing derived fpqc/base-change and blow-up-construction prerequisites remain named references; later additive-invariant and cdh results are unselected context. All proof text and formulas below are editable author TeX. Mathematical wording, language and proof order are retained. Only the article class, title matter and bibliography presentation are adapted for independent compilation; no proof has been generated or repaired.
\par\noindent\textbf{Editorial difficulty:} H3: The authored argument builds derived virtual-Cartier duality, combines it with derived projective-bundle semi-orthogonality, and proves generation and preservation of perfect complexes under quasi-smooth blow-up functors using fpqc descent. This requires advanced derived algebraic geometry and stable infinity-categories beyond ordinary qualifying exams.
\par\noindent\textbf{Editorial retrieval and attribution:} Retrieved 2026-10-01. DOI 10.5802/ahl.55. Source attribution notice: ../licenses/2634-source-license.txt; complete legal text: ../licenses/CC-BY-4.0.txt. Original source excerpts and the existing source-material check are retained in ../sources/2634/.
\medskip\hrule\medskip
\noindent\textbf{Author statement, complete proof and local supporting text follow in source order.}
\section{Introduction}\label{sec:intro}

Let $X$ be a scheme and $i : Z \to X$ a regular closed immersion. This means that $Z$ is, Zariski-locally on $X$, the zero-locus of some regular sequence of functions $f_1,\,\ldots,\,f_n \in \Gamma(X,\sO_X)$. Then the blow-up $\Bl_{Z/X}$ fits into a square
\begin{equation}\label{eq:blow-up square}
\begin{tikzcd}
\P\left(\sN_{Z/X}\right) \ar{r}{i_D}\ar{d}{q} & \Bl_{Z/X} \ar{d}{p}
\\
Z \ar{r}{i} & X,
\end{tikzcd}
\end{equation}
where the exceptional divisor is the projective bundle associated to the conormal sheaf $\sN_{Z/X}$, which under the assumptions is locally free of rank $n$. A result of Thomason~\cite{ThomasonBlowup} asserts that after taking algebraic K-theory, the induced square of spectra
\[
\begin{tikzcd}
\K(X) \ar{r}{i^*}\ar{d}{p^*} & \K(Z) \ar{d}
\\
\K\left(\Bl_{Z/X}\right) \ar{r} & \K\left(\P\left(\sN_{Z/X}\right)\right)
\end{tikzcd}
\]
is homotopy cartesian. Here $\K(X)$ denotes the Bass--Thomason--Trobaugh algebraic K-theory spectrum of perfect complexes on a scheme $X$. We may summarize this property by saying that algebraic K-theory satisfies \emph{descent} with respect to blow-ups in regularly immersed centres.

Now suppose that $i$ is more generally a \emph{quasi-smooth} closed immersion of derived schemes. This means that $Z$ is, Zariski-locally on $X$, the \emph{derived} zero-locus of some arbitrary sequence of functions $f_1,\,\ldots,\,f_n \in \Gamma(X,\sO_X)$. (When $X$ is a classical scheme and the sequence is regular, this is the same as the classical zero-locus, and we are in the situation discussed above.) In the derived setting there is still a conormal sheaf $\sN_{Z/X}$ on $Z$, locally free of rank $n$, and one may still form the blow-up square~\eqref{eq:blow-up square}, see~\cite{KhanBlowup}. Our goal in this paper is to generalize Thomason's result above to this situation. At the same time we also allow $X$ to be a derived \emph{Artin stack}, and consider any \emph{additive invariant} of stable $\infty$-categories (see
Definition~\ref{defn:additive invariant}). Examples of additive invariants include algebraic K-theory $\K$, connective algebraic K-theory $\K^\mrm{cn}$, topological Hochschild homology $\mrm{THH}$, and topological cyclic homology $\mrm{TC}$.

\begingroup
\begin{thmX}\label{thm:blow-up descent}
Let $E$ be an additive invariant of stable $\infty$-categories. Then $E$ satisfies descent by quasi-smooth blow-ups. That is, given a derived Artin stack $X$ and a quasi-smooth closed immersion $i : Z \to X$ of virtual codimension $n\ge 1$, form the blow-up square~\eqref{eq:blow-up square}. Then the induced commutative square
\[
\begin{tikzcd}
E(X) \ar{r}{i^*}\ar{d}{p^*} & E(Z) \ar{d}
\\
E\left(\Bl_{Z/X}\right) \ar{r} & E\left(\P\left(\sN_{Z/X}\right)\right)
\end{tikzcd}
\]
is homotopy cartesian.
\end{thmX}
\endgroup

We deduce Theorem~\ref{thm:blow-up descent} from an analysis of the categories of perfect complexes on $\Bl_{Z/X}$ and on the exceptional divisor $\P(\sN_{Z/X})$. The relevant notion is that of a \emph{semi-orthogonal decomposition}, see
Definition~\ref{defn:SOD}.
%%Theorem A à Cintroduire un setcounter spécifique pour les trois enonces.

\begin{thmX}\label{thm:Perf(P(E))}
Let $X$ be a derived Artin stack. For any locally free $\sO_X$-module $\sE$ of rank $n+1$, $n\ge 0$, consider the projective bundle $q : \P(\sE) \to X$. Then we have:
\begin{enumerate}\romanenumi
\item\label{item:Perf(P(E))/fully faithful}
For each $0\le k\le n$, the assignment $\sF \mapsto q^*(\sF) \otimes \sO(-k)$ defines a fully faithful functor $\Perf(X) \to \Perf(\P(\sE))$, whose essential image we denote $\bA(-k)$.

\item\label{item:Perf(P(E))/SOD}
The sequence of full subcategories $(\bA(0),\ldots,\bA(-n))$ forms a semi-orthogonal decomposition of $\Perf(\P(\sE))$.
\end{enumerate}
\end{thmX}

\begin{thmX}\label{thm:Perf(Bl)}
Let $X$ be a derived Artin stack. For any quasi-smooth closed immersion $i : Z \to X$ of virtual codimension $n\ge 1$, form the blow-up square~\eqref{eq:blow-up square}. Then we have:
\begin{enumerate}\romanenumi
\item\label{item:Perf(Bl)/fully faithful 1}
The assignment $\sF \mapsto p^*(\sF)$ defines a fully faithful functor $\Perf(X) \linebreak\to \Perf(\Bl_{Z/X})$, whose essential image we denote $\bB(0)$.

\item\label{item:Perf(Bl)/fully faithful 2}
For each $1\le k\le n-1$, the assignment $\sF \mapsto (i_D)_*(q^*(\sF) \otimes \sO(-k))$ defines a fully faithful functor $\Perf(Z) \to \Perf(\Bl_{Z/X})$, whose essential image we denote $\bB(-k)$.

\item\label{item:Perf(Bl)/SOD}
The sequence of full subcategories $(\bB(0),\,\ldots,\,\bB(-n+1))$ forms a \emph{semi-orthogonal decomposition} of $\Perf(\Bl_{Z/X})$.
\end{enumerate}
\end{thmX}

We immediately deduce the projective bundle and blow-up formulas
\[
E(\P(\sE)) \simeq \bigoplus_{m=0}^{n} E(X),
\quad E\left(\Bl_{Z/X}\right) \simeq E(X) \oplus \bigoplus_{k=1}^{n-1} E(Z),
\]
for any additive invariant $E$, see Corollaries~\ref{cor:E(P(E))} and~\ref{cor:E(Bl_{Z/X})}, from which Theorem~\ref{thm:blow-up descent} immediately follows (see Subsection~\ref{ssec:blowup/additive}).

The results mentioned above admit the following interesting special cases:
\begin{enumerate}
\item Suppose that $X$ is a smooth projective variety over the field of complex numbers. This case of Theorem~\ref{thm:Perf(P(E))} was proven by Orlov in~\cite{OrlovSOD}. He also proved Theorem~\ref{thm:Perf(Bl)} for any \emph{smooth} subvariety $Z \hook X$.
\item More generally suppose that $X$ is a quasi-compact quasi-separated classical scheme. Then the projective bundle formula (Corollary~\ref{cor:E(P(E))}) for algebraic K-theory was proven by Thomason~\cite{ThomasonProjectiveBundle,ThomasonTrobaugh}. Similarly suppose that $i : Z \to X$ is a quasi-smooth closed immersion of quasi-compact quasi-separated classical schemes. Then it is automatically a regular closed immersion, and in this case Thomason also proved Corollary~\ref{cor:E(Bl_{Z/X})} for algebraic K-theory~\cite{ThomasonBlowup}. In fact, the papers~\cite{ThomasonProjectiveBundle} and~\cite{ThomasonBlowup} essentially contain under these assumptions proofs of Theorems~\ref{thm:Perf(P(E))} and~\ref{thm:Perf(Bl)}, respectively, even if the term ``semi-orthogonal decomposition'' is not used explicitly. For $\mrm{THH}$ and $\mrm{TC}$, these cases of Corollaries~\ref{cor:E(P(E))} and~\ref{cor:E(Bl_{Z/X})} were proven by Blumberg and Mandell~\cite{BlumbergMandellTHH}.

\item More generally still, let $X$ and $Z$ be classical Artin stacks. These cases of Theorems~\ref{thm:Perf(P(E))} and~\ref{thm:Perf(Bl)} are proven by Bergh and Schn\"urer in~\cite{BerghSchnuerer}. However we note that Corollaries~\ref{cor:E(P(E))} and~\ref{cor:E(Bl_{Z/X})} were obtained earlier by Krishna and Ravi in~\cite{KrishnaRavi}, and their arguments in fact prove Theorems~\ref{thm:Perf(P(E))} and~\ref{thm:Perf(Bl)} for classical Artin stacks.

\item Let $X$ be a noetherian affine classical scheme, and let $Z$ be the derived zero-locus of some functions $f_1,\,\ldots,\,f_n \in \Gamma(X, \sO_X)$. Then the canonical morphism $i : Z \to X$ is a quasi-smooth closed immersion. In this case, Theorem~\ref{thm:blow-up descent} for algebraic K-theory was proven by Kerz--Strunk--Tamme~\cite{KerzStrunkTamme} (where the blow-up $\Bl_{Z/X}$ was explicitly modelled as the derived fibred product $X \fibprod_{\A^n} \Bl_{\{0\}/\A^n}$), as part of their proof of Weibel's conjecture on negative K-theory.
\end{enumerate}

Let $\KH$ denote homotopy invariant K-theory. Recall that this is the $\A^1$-localization of the presheaf $X \mapsto \K(X)$. That is, it is obtained by forcing the property of $\A^1$-homotopy invariance: for every quasi-compact quasi-separated algebraic space $X$, the map
\[
\KH(X) \to \KH\left(X \times \A^1\right)
\]
is invertible (see~\cite{CisinskiKH,WeibelKH}). As an application of Theorem~\ref{thm:blow-up descent}, we give a new proof of the following theorem of Cisinski~\cite{CisinskiKH}:

\begin{thmX}\label{thm:KH cdh descent}
The presheaf of spectra $S \mapsto \KH(S)$ satisfies cdh descent on the site of quasi-compact quasi-separated algebraic spaces.
\end{thmX}

This was first proven by Haesemeyer~\cite{Haesemeyer} for schemes over a field of characteristic zero, using resolution of singularities. Cisinski's proof over general bases (noetherian schemes of finite dimension) relies on Ayoub's proper base change theorem in motivic homotopy theory. A different proof of Theorem~\ref{thm:KH cdh descent} (also in the noetherian setting) was recently given by Kerz--Strunk--Tamme~\cite[Theorem~C]{KerzStrunkTamme}, as an application of pro-cdh descent and their resolution of Weibel's conjecture on negative K-theory. The proof we give here is more direct and uses a new criterion for cdh descent (see Theorem~\ref{thm:cdh criterion} for a more precise statement):

\begin{thmX}\label{thm:cdh intro}
Let $\sF$ be a Nisnevich sheaf of spectra on the category of quasi-compact quasi-separated algebraic spaces. Then $\sF$ satisfies cdh descent if and only if it sends closed squares and quasi-smooth blow-up squares to cartesian squares.
\end{thmX}

Theorem~\ref{thm:cdh intro} can be compared to a similar criterion due to Haesemeyer, implicit in~\cite{Haesemeyer}, which applies to Nisnevich sheaves of spectra on the category of schemes over a field $k$ of characteristic zero. It asserts that for such a sheaf, cdh descent is equivalent to descent for finite cdh squares and regularly immersed blow-up squares. Note that the first condition is stronger than descent for closed squares, while the second is weaker than descent for quasi-smooth blow-up squares: regularly immersed blow-up squares are precisely those quasi-smooth blow-up squares where all schemes appearing are underived. For invariants of stable $\infty$-categories, a similar cdh descent criterion was noticed independently by Land and Tamme~\cite[Theorem.~A.2]{LandTamme}. The h-descent criterion of Halpern-Leistner and Preygel~\cite[Proposition~3.2.1]{HalpernLeistnerPreygel} is also somewhat similar in spirit.

Theorem~\ref{thm:KH cdh descent} was extended to certain nice Artin stacks recently by Hoyois and Hoyois--Krishna~\cite{HoyoisCdh,HoyoisKrishna}. Our cdh descent criterion also applies in that setting (Remark~\ref{rem:variants of cdh criterion}$\MK$\eqref{item:variants of cdh criterion/stacks}) and gives another potential approach to such results.

The organization of this paper is as follows. We begin in Section~\ref{sec:prelim} with some background on derived algebraic geometry and on semi-orthogonal decompositions\linebreak of stable $\infty$-categories.
Section~\ref{sec:proj} is dedicated to the proof of Theorem~\ref{thm:Perf(P(E))}.\linebreak We first show that the semi-orthogonal decomposition exists on the larger stable $\infty$-category $\Qcoh(\P(\sE))$ (Theorem~\ref{thm:Qcoh(P(E))}). Then we show that it restricts to $\Perf(\P(\sE))$ (Subsection~\ref{ssec:proof of Perf(P(E))}), and deduce the projective bundle formula (Corollary~\ref{cor:E(P(E))}) for any additive invariant.

We follow a similar pattern in Section~\ref{sec:blowup} to prove Theorem~\ref{thm:Perf(Bl)}. There is a semi-orthogonal decomposition on $\Qcoh(\Bl_{Z/X})$ (Theorem~\ref{thm:Qcoh(Bl)}) which then restricts to $\Perf(\Bl_{Z/X})$ (Subsection~\ref{ssec:proof of Perf(Bl)}). This gives both the blow-up formula (Corollary~\ref{cor:E(Bl_{Z/X})}) as well as Theorem~\ref{thm:blow-up descent} (\ref{sssec:proof of blow-up descent}) for additive invariants. As input we prove a Grothendieck duality statement for virtual Cartier divisors (Proposition~\ref{prop:omega_D/X}) that should be of independent interest.

Section~\ref{sec:cdh} contains our results on cdh descent and KH. We first give the general cdh descent criterion (Theorem~\ref{thm:cdh criterion}). We apply this criterion to $\KH$ to give our proof of Theorem~\ref{thm:KH cdh descent} (\ref{sssec:proof of KH cdh descent}).

I would like to thank Marc Hoyois, Charanya Ravi, and David Rydh for helpful discussions and comments on previous revisions. I am especially grateful to David Rydh for pointing out the relevance of the resolution property in Section~\ref{sec:cdh}.


\section{Preliminaries}\label{sec:prelim}

Throughout the paper we work with the language of $\infty$-categories as in~\cite{HTT,HA-20170918}.

\subsection{Derived algebraic geometry}

This paper is set in the world of derived algebraic geometry, as in~\cite{GaitsgoryRozenblyum,SAG-20180204,HAG2}.

\subsubsection{Derived schemes and stacks}
Let $\SCRing$ denote the $\infty$-category of simplicial commutative rings. A \emph{derived stack} is an étale sheaf of spaces $X : \SCRing \to \Spc$. If $X$ is corepresentable by a simplicial commutative ring $A$, we write $X = \Spec(A)$ and call $X$ an \emph{affine derived scheme}. A \emph{derived scheme} is a derived stack $X$ that admits a Zariski atlas by affine derived schemes, i.e., a jointly surjective family $(U_i \to X)_i$ of Zariski open immersions with each $U_i$ an affine derived scheme. Allowing Nisnevich, étale or smooth atlases, respectively, gives rise to the notions of \emph{derived algebraic space}\footnote{That this agrees with the classical notion of algebraic space (at least under quasi-compactness and quasi-separatedness hypotheses) follows from~\cite[Proposition~5.7.6]{RaynaudGruson}. That it agrees with Lurie's definition follows from~\cite[Example~3.7.1.5]{SAG-20180204}.} \emph{derived Deligne--Mumford stack}, and \emph{derived Artin stack}.
 The precise definition is slightly more involved, see e.g.~\cite[Volume~I, Section~4.1]{GaitsgoryRozenblyum}.

Any derived stack $X$ admits an underlying classical stack which we denote $X_\cl$. If $X$ is a derived scheme, algebraic space, Deligne--Mumford or Artin stack, then $X_\cl$ is a classical such. For example, $\Spec(A)_\cl = \Spec(\pi_0(A))$ for a simplicial commutative ring $A$.

\subsubsection{Quasi-smooth morphisms}
Let $X$ be a derived scheme and let $f_1,\,\ldots,\,f_n \in \Gamma(X, \sO_X)$ be functions classifying a morphism $f : X \to \A^n$ to affine space. The \emph{derived zero-locus} of these functions is given by the derived fibred product
\begin{equation}\label{eq:quasi-smooth}
\begin{tikzcd}
Z \ar{r}\ar{d} & X \ar{d}{f}
\\
\{0\} \ar{r} & \A^n.
\end{tikzcd}
\end{equation}
If $X$ is classical, then $Z$ is classical if and only if the sequence $(f_1,\,\ldots,\,f_n)$ is regular in the sense of~\cite{SGA6}, in which case $Z$ is regularly immersed. A closed immersion of derived schemes $i : Z \to X$ is called \emph{quasi-smooth} (of virtual codimension $n$) if it is cut out Zariski-locally as the derived zero-locus of $n$ functions on $X$. Equivalently, this means that $i$ is of finite presentation and its shifted cotangent complex\linebreak$\sN_{Z/X}:= \sL_{Z/X}[-1]$ is locally free (of rank $n$). A closed immersion of derived Artin stacks is quasi-smooth if it satisfies this condition smooth-locally.

A morphism of derived schemes $f : Y \to X$ is quasi-smooth if it can be factored, Zariski-locally on $Y$, through a quasi-smooth closed immersion $i : Y \to X'$ and a smooth morphism $X' \to X$. A morphism of derived Artin stacks is quasi-smooth if it satisfies this condition smooth-locally on $Y$. We refer to~\cite{KhanBlowup} for more details on quasi-smoothness.

\subsubsection{Blow-ups}

Important for us is the following construction from~\cite{KhanBlowup}. Given any quasi-smooth closed immersion $i : Z \to X$ of derived Artin stacks, there is an associated \emph{quasi-smooth blow-up square}:
\begin{equation}\label{eq:qsbl square}
\begin{tikzcd}
D \ar{r}{i_D}\ar{d}{q} & \Bl_{Z/X} \ar{d}{p}
\\
Z \ar{r}{i} & X.
\end{tikzcd}
\end{equation}
Here $\Bl_{Z/X}$ is the blow-up of $X$ in $Z$, which is a quasi-smooth proper derived Artin stack over $X$, and $D = \P(\sN_{Z/X})$ is the projectivized normal bundle, which is a smooth proper derived Artin stack over $Z$. This square is universal with the following properties:
\begin{enumerate}\alphenumi
\item the morphism $i_D$ is a quasi-smooth closed immersion of virtual codimension~$1$, i.e., a virtual effective Cartier divisor;
\item the underlying square of classical Artin stacks is cartesian; and
\item the canonical map $q^*\sN_{Z/X} \to \sN_{D/\Bl_{Z/X}}$ is surjective on $\pi_0$. When $X$ is a derived scheme (resp.~derived algebraic space, derived Deligne--Mumford stack), then so is $\Bl_{Z/X}$.
\end{enumerate}

\subsubsection{Quasi-coherent sheaves}

Given a derived stack $X$, the stable $\infty$-category of quasi-coherent sheaves $\Qcoh(X)$ is the limit
\[
\Qcoh(X) = \lim_{\Spec(A) \to X} \Qcoh(\Spec(A))
\]
taken over all morphisms $\Spec(A) \to X$ with $A \in \SCRing$. Here $\Qcoh(\Spec(A))$ is the stable $\infty$-category $\Mod_A$ of $A$-modules\footnotemark~in the sense of Lurie.
\footnotetext{Note that if $A$ is discrete (an ordinary commutative ring), then this is not the abelian category of discrete $A$-modules, but rather the derived $\infty$-category of this abelian category as in~\cite[Chapter~1]{HA-20170918}.} Informally speaking, a quasi-coherent sheaf $\sF$ on $X$ is thus a collection of quasi-coherent sheaves $x^*(\sF) \in \Qcoh(\Spec(A))$, for every simplicial commutative ring $A$ and every $A$-point $x : \Spec(A) \to X$, together with a homotopy coherent system of compatibilities.

The full subcategory $\Perf(X) \subset \Qcoh(X)$ is similarly the limit
\[
\Perf(X) = \lim_{\Spec(A) \to X} \Perf(\Spec(A)),
\]
where $\Perf(\Spec(A))$ is the stable $\infty$-category $\Mod_A^\perf$ of \emph{perfect} $A$-modules. In other words, $\sF \in \Qcoh(X)$ belongs to $\Perf(X)$ if and only if $x^*(\sF)$ is perfect for every simplicial commutative ring $A$ and every morphism $x : \Spec(A) \to X$.

\subsubsection{Functoriality}\label{sssec:Qcoh functoriality}

There is an inverse image functor $f^* : \Qcoh(X) \to \Qcoh(Y)$ for any morphism of derived stacks $f : Y \to X$. It preserves perfect complexes and induces a functor $f^* : \Perf(X) \to \Perf(Y)$. Regarded as presheaves of $\infty$-categories, the assignments $X \mapsto \Qcoh(X)$ and $X \mapsto \Perf(X)$ satisfy \emph{descent} for the fpqc topology (\cite[Corollary~D.6.3.3]{SAG-20180204}, \cite[Theorem~1.3.4]{GaitsgoryRozenblyum}). This means in particular that given any fpqc covering family $(f_\alpha : X_\alpha \to X)_\alpha$, the family of inverse image functors $f_\alpha^* : \Qcoh(X) \to \Qcoh(X_\alpha)$ is jointly conservative.

If $f : Y \to X$ is quasi-compact and \emph{schematic}, in the sense that its fibre over any affine derived scheme is a derived scheme, then there is a direct image functor $f_*$, right adjoint to $f^*$, which commutes with colimits and satisfies a base change formula against inverse images (\cite[Proposition~2.5.4.5]{SAG-20180204}, \cite[Volume~1, Chapter~3, Proposition~2.2.2]{GaitsgoryRozenblyum}). If $f$ is proper, locally of finite presentation, and of finite tor-amplitude, then $f_*$ also preserves perfect complexes~\cite[Theorem~6.1.3.2]{SAG-20180204}.

\subsection{Semi-orthogonal decompositions}

The following definitions were originally formulated by~\cite{BondalKapranov} in the language of triangulated categories and are standard.

\begin{defi}
Let $\bC$ be a stable $\infty$-category and $\bD$ a stable full subcategory. An object $x \in \bC$ is \emph{left orthogonal}, resp.~\emph{right orthogonal}, to $\bD$ if the mapping space $\Maps_\bC(x,d)$, resp.~$\Maps_\bC(d,x)$, is contractible for all objects $d\in\bD$. We let $^\perp\bD \subseteq \bC$ and $\bD^\perp \subseteq \bC$ denote the full subcategories of left orthogonal and right orthogonal objects, respectively.
\end{defi}

\begin{defi}\label{defn:SOD}
Let $\bC$ be a stable $\infty$-category and let $\bC(0),\,\ldots,\,\bC(-n)$ be full stable subcategories. Suppose that the following conditions hold:
\begin{enumerate}\romanenumi
\item For all integers $i>j$, there is an inclusion $\bC(i) \subseteq I\,^\perp\bC(j)$.
\item The $\infty$-category $\bC$ is generated by the subcategories $\bC(0),\,\ldots,\,\bC(-n)$, under finite limits and finite colimits.
\end{enumerate}
Then we say that the sequence $(\bC(0),\,\ldots,\,\bC(-n))$ forms a \emph{semi-orthogonal decomposition} of $\bC$.
\end{defi}

Semi-orthogonal decompositions of length $2$ come from \emph{split short exact sequences} of stable $\infty$-categories, as in~\cite{BlumbergGepnerTabuada}.

\begin{defi} \ \\*[-1.5em]%\leavevmode
\begin{enumerate}\romanenumi
\item A \emph{short exact sequence} of small stable $\infty$-categories is a diagram
\[
\bC' \xrightarrow{i} \bC \xrightarrow{p} \bC'',
\]
where $i$ and $p$ are exact, the composite $p\circ i$ is null-homotopic, $i$ is fully faithful, and $p$ induces an equivalence $(\bC/\bC')^\idem \simeq (\bC'')^\idem$ (where $(-)^\idem$ denotes idempotent completion).
\item A short exact sequence of small stable $\infty$-categories
\[
\bC' \xrightarrow{i} \bC \xrightarrow{p} \bC''
\]
is \emph{split} if there exist functors $q : \bC \to \bC'$ and $j : \bC'' \to \bC$, right adjoint to $i$ and $p$, respectively, such that the unit $\id \to q\circ i$ and co-unit $p\circ j \to \id$ are invertible.
\end{enumerate}
\end{defi}

\begin{rema}\label{rem:SOD to split SES}
Let $\bC$ be a small stable $\infty$-category, and let $(\bC(0), \bC(-1))$ be a semi-orthogonal decomposition. Then for any object $x \in \bC$, there exists an exact triangle
\[
x(0) \to x \to x(-1),
\]
where $x(0) \in \bC(0)$ and $x(-1) \in \bC(-1)$. To see this, simply observe that the full subcategory spanned by objects $x$ for which such a triangle exists, is closed under finite limits and colimits, and contains $\bC(0)$ and $\bC(-1)$. Moreover, the assignments $x \mapsto x(0)$ and $x \mapsto x(-1)$ determine well-defined functors $q : \bC \to \bC(0)$ and $p : \bC \to \bC(-1)$, respectively, which are right and left adjoint, respectively, to the inclusions (see e.g.~\cite[Remark~7.2.0.2]{SAG-20180204}). It follows from this that any semi-orthogonal decomposition $(\bC(0),\bC(-1))$ induces a split short exact sequence
\[
\bC(0) \to \bC \xrightarrow{p} \bC(-1).
\]
\end{rema}

\begin{lemm}\label{lem:SOD filtration}
Let $\bC$ be a stable $\infty$-category, and let $(\bC(0),\,\ldots,\,\bC(-n))$ be a sequence of full stable subcategories forming a semi-orthogonal decomposition of $\bC$. For each $0 \le m \le n$, let $\bC_{\le -m} \subseteq \bC$ denote the full stable subcategory generated by objects in the union $\bC(-m)\cup\,\cdots\,\cup\bC(-n)$, and let $\bC_{\le -n-1} \subseteq \bC$ denote the full subcategory spanned by the zero object. Then there are split short exact sequences
\[
\bC(-m) \hook \bC_{\le -m} \to \bC_{\le -m-1}
\]
for each $0\le m\le n$.
\end{lemm}

\begin{proof}
It follows from the definitions that for each $0\le m\le n$, the sequence $(\bC(-m), \bC_{\le -m-1})$ forms a semi-orthogonal decomposition of $\bC_{\le -m}$. Therefore the claim follows from Remark~\ref{rem:SOD to split SES}.
\end{proof}

\subsection{Additive and localizing invariants}

The following definition is from~\cite{BlumbergGepnerTabuada}, except that we do not require commutativity with filtered colimits.

\begin{defi}\label{defn:additive invariant}
Let $\bA$ be a stable presentable $\infty$-category. Let $E$ be an $\bA$-valued functor from the $\infty$-category of small stable $\infty$-categories and exact functors.

\begin{enumerate}\romanenumi
\item We say that $E$ is an \emph{additive invariant} if for any split short exact sequence
\[
\bC' \xrightarrow{i} \bC \xrightarrow{p} \bC'',
\]
the induced map
\[
E(\bC') \oplus E(\bC'') \xrightarrow{(i,j)} E(\bC)
\]
is invertible, where $j$ is a right adjoint to $p$.
\item We say that $E$ is a \emph{localizing invariant} if for any short exact sequence
\[
\bC' \xrightarrow{i} \bC \xrightarrow{p} \bC'',
\]
the induced diagram
\[
E(\bC') \to E(\bC) \to E(\bC'')
\]
is an exact triangle.
\end{enumerate}
\end{defi}

\begin{rema}
Any localizing invariant is also additive.
\end{rema}

\begin{lemm}\label{lem:additive}
Let $\bC$ be a stable $\infty$-category, and let $(\bC(0),\,\ldots,\,\bC(-n))$ be a sequence of full stable subcategories forming a semi-orthogonal decomposition of $\bC$. Then for any additive invariant $E$ there is a canonical isomorphism
\[
E(\bC) \simeq \bigoplus_{m=0}^n E(\bC(-m)).
\]
\end{lemm}

\begin{proof}
Follows immediately from Lemma~\ref{lem:SOD filtration}.
\end{proof}

\section{The projective bundle formula}\label{sec:proj}

\subsection{Projective bundles}\label{ssec:proj/proj}

Let $X$ be a derived stack and $\sE$ a locally free $\sO_X$-module of finite rank. Recall that the \emph{projective bundle} associated to $\sE$ is a derived stack $\P(\sE)$ over $X$ equipped with an invertible sheaf $\sO(1)$ together with a surjection $\sE \to \sO(1)$. More precisely, for any derived scheme $S$ over $X$, with structural morphism $x : S \to X$, the space of $S$-points of $\P(\sE)$ is the space of pairs $(\sL, u)$, where $\sL$ is a locally free $\sO_S$-module of rank $1$, and $u : x^*(\sE) \to \sL$ is surjective on $\pi_0$. We recall the standard properties of this construction:

\begin{prop} \ \\*[-1.5em]
\begin{enumerate}\romanenumi
\item If $f : X' \to X$ is a morphism of derived stacks, then there is a canonical isomorphism $\P(f^*(\sE)) \to \P(\sE) \fibprod_X X'$ of derived stacks over $X'$.
\item The projection $\P(\sE)\to X$ is proper and schematic. In particular, if $X$ is a derived scheme (resp.~derived algebraic space, derived Deligne--Mumford stack, derived Artin stack), then the same holds for the derived stack $\P(\sE)$.
\item The relative cotangent complex $\sL_{\P(\sE)/X}$ is canonically isomorphic to $\sF \otimes \sO(-1)$, where the locally free sheaf $\sF$ is the fibre of the canonical map $\sE \to \sO(1)$. In particular, the morphism $\P(\sE) \to X$ is smooth of relative dimension equal to $\rk(\sE) - 1$.
\end{enumerate}
\end{prop}

\begin{prop}[Serre]\label{prop:Serre}
Let $X$ be a derived Artin stack, and $\sE$ a locally free sheaf of rank $n+1$, $n\ge 0$. If $q : \P(\sE) \to X$ denotes the associated projective bundle, then we have canonical isomorphisms
\[
q_*(\sO(0)) \simeq \sO_X,\quad q_*(\sO(-m)) \simeq 0 \; (1\le m\le n)
\]
in $\Qcoh(X)$.
\end{prop}

\begin{proof}
There is a canonical map $\sO_X \to q_*(\sO(0))$, the unit of the adjunction $(q^*,q_*)$, and there is a unique map $0 \to q_*(\sO(-m))$ for each $m$. To show that these are invertible, we may use fpqc descent and base change to the case where $X$ is affine and $\sE$ is free. Then this is Serre's computation, as generalized to the derived setting by Lurie~\cite[Theorem.~5.4.2.6]{SAG-20180204}.
\end{proof}

\subsection{Semi-orthogonal decomposition on \texorpdfstring{$\Qcoh(\P(\sE))$}{Qcoh(P(E))}}\label{ssec:proj/qcoh}

In this subsection we will show that the stable $\infty$-category $\Qcoh(\P(\sE))$ admits a canonical semi-orthogonal decomposition.

\begin{theo}\label{thm:Qcoh(P(E))}
Let $X$ be a derived Artin stack. Let $\sE$ be a locally free $\sO_X$-module of rank $n+1$, $n\ge 0$, and $q : \P(\sE) \to X$ the associated projective bundle. Then we have:
\begin{enumerate}\romanenumi
\item\label{item:Qcoh(P(E))/fully faithful}
For every integer $k\in\bZ$, the assignment $\sF \mapsto q^*(\sF)\otimes\sO(k)$ defines a fully faithful functor $\Qcoh(X) \to \Qcoh(\P(\sE))$.
\item\label{item:Qcoh(P(E))/SOD}
For every integer $k\in\bZ$, let $\bC(k) \subset \Qcoh(\P(\sE))$ denote the essential image of the functor in (i). Then the subcategories $\bC(k),\,\ldots,\,\bC(k-n)$ form a semi-orthogonal decomposition of $\Qcoh(\P(\sE))$.
\end{enumerate}
\end{theo}

We will need the following facts (see~\cite[Lemmas~7.2.2.2 and 5.6.2.2]{SAG-20180204}):

\begin{lemm}\label{lem:O(n+1)}
Let $R$ be a simplicial commutative ring and $X = \Spec(R)$. Denote by $\P^n_R = \P(\sO^{n+1}_X)$ the $n$-dimensional projective space over $R$. Then for every integer $m\in\bZ$, there is a canonical isomorphism
\[
\colim_{J \subsetneq [n]} \sO(m+\abs{J}) \isoto \sO(m+n+1)
\]
in $\Qcoh(\P^n_R)$, where the colimit is taken over the proper subsets $J$ of the set\linebreak$[n] = \{0,1,\,\ldots,\,n\}$, and $0\le\abs{J}\le n$ denotes the cardinality of such a subset.
\end{lemm}

\begin{lemm}\label{lem:surjection from O(m)'s}
Let $R$ be a simplicial commutative ring and $X = \Spec(R)$. Denote by $\P^n_R = \P(\sO^{n+1}_X)$ the $n$-dimensional projective space over $R$. Then for any connective quasi-coherent sheaf $\sF \in \Qcoh(\P^n_R)$, there exists a map
\[
\bigoplus_\alpha \sO(d_\alpha) \to \sF,
\]
with $d_\alpha \in \bZ$, which is surjective on $\pi_0$.
\end{lemm}

\begin{proof}[Proof of Theorem~\ref{thm:Qcoh(P(E))}]
Since the functors $-\otimes\sO(k)$ are equivalences, it will suffice to take $k=0$ in both claims. For claim~\eqref{item:Qcoh(P(E))/fully faithful} we want to show that the unit map $\sF \to q_*q^*(\sF)$ is invertible for all $\sF \in \Qcoh(X)$. By fpqc descent and base change (Subsubsection~\ref{sssec:Qcoh functoriality}), we may reduce to the case where $X = \Spec(R)$ is affine and $\sE = \sO^{n+1}_S$ is free. Now both functors $q^*$ and $q_*$ are exact and moreover commute with arbitrary colimits (the latter by Subsubsection~\ref{sssec:Qcoh functoriality} since $q$ is quasi-compact and schematic), and $\Qcoh(X) \simeq \Mod_R$ is generated by $\sO_X$ under colimits and finite limits. Therefore we may assume $\sF = \sO_X$, in which case the claim holds by Proposition~\ref{prop:Serre}.

For claim~\eqref{item:Qcoh(P(E))/SOD}, let us first check the orthogonality condition in Definition~\ref{defn:SOD}. Thus take $\sF,\sG\in \Qcoh(X)$ and consider the mapping space
\[
\Maps(q^*(\sF), q^*(\sG) \otimes \sO(-m)) \simeq \Maps(\sF, q_*(\sO(-m)) \otimes \sG)
\]
for $1\le m\le n$, where the identification results from the projection formula. Since $q_*(\sO(-m)) \simeq 0$ by Proposition~\ref{prop:Serre}, this space is contractible.

It now remains to show that every $\sF \in \Qcoh(\P(\sE))$ belongs to the full subcategory $\langle \bC(0),\,\ldots,\,\bC(-n) \rangle \subseteq \Qcoh(\P(\sE))$ generated under finite colimits and limits by the subcategories $\bC(0),\,\ldots,\,\bC(-n)$. Set $\sG_{-1} = \sF \otimes \sO(-1)$ and define $\sG_m$, for $m\ge 0$, so that we have exact triangles
\begin{equation}\label{eq:G_m+1}
q^*q_*(\sG_{m-1} \otimes \sO(1)) \xrightarrow{\mrm{counit}} \sG_{m-1} \otimes \sO(1) \to \sG_{m}.
\end{equation}
For each $m\ge -1$, we claim that $\sG_m$ is right orthogonal to each of the subcategories $\bC(0),\bC(1),\,\ldots,\,\bC(m)$. For $m=-1$ the claim is vacuous, so take $m\ge 0$ and assume by induction that it holds for $m-1$. Since $q^*q_*(\sG_{m-1} \otimes \sO(1))$ is contained in $\bC(0)$, it follows that $\sG_{m}$ is right orthogonal to $\bC(0)$. To show that $\sG_m$ is right orthogonal to $\bC(i)$, for $1\le i\le m$, it will suffice to show that the left-hand and middle terms of the exact triangle~\eqref{eq:G_m+1} are both right orthogonal to $\bC(i)$. For the left-hand term this follows from the inclusion $\bC(0) \subset \bC(i)^\perp$, demonstrated above. For the middle term $\sG_{m-1} \otimes \sO(1)$, the claim follows by the induction hypothesis.

Now we claim that $\sG_n$ is zero. Using fpqc descent again, we may assume that $X = \Spec(R)$ and $\sE = \sO^{\oplus n+1}_X$ is free (since the sequence $(\sG_{-1},\sG_0,\,\ldots,\,\sG_n)$ is stable under base change). Using Lemma~\ref{lem:surjection from O(m)'s} we can build a map
\[
\varphi : \bigoplus_\alpha \sO(m_\alpha)[k_\alpha] \to \sG_n
\]
which is surjective on all homotopy groups. From Lemma~\ref{lem:O(n+1)} it follows that $\sG_n$ is right orthogonal to all $\bC(i)$, $i\in\bZ$. Thus $\varphi$ must be null-homotopic, so $\sG_n \simeq 0$ as claimed. Working backwards, we deduce that $\sG_{n-1} \in \bC(-1)$,\dots, $\sG_0 \in \langle \bC(-1),\,\ldots,\,\bC(-n)\rangle$, and then finally that $\sF \in \langle \bC(0), \bC(-1),\,\ldots,\,\bC(-n)\rangle$ as claimed.
\end{proof}

\subsection{Proof of Theorem~\ref{thm:Perf(P(E))}}\label{ssec:proof of Perf(P(E))}

We now deduce Theorem~\ref{thm:Perf(P(E))} from Theorem~\ref{thm:Qcoh(P(E))}. First note that the fully faithful functor $\sF \mapsto q^*(\sF)\otimes\sO(k)$ of Theorem~\ref{thm:Qcoh(P(E))}$\MK$\eqref{item:Qcoh(P(E))/fully faithful} restricts to a fully faithful functor $\Perf(X) \to \Perf(\P(\sE))$, since $q^*$ preserves perfect complexes. This shows Theorem~\ref{thm:Perf(P(E))}$\MK$\eqref{item:Perf(P(E))/fully faithful}.

For part~\eqref{item:Perf(P(E))/SOD} we argue again as in the proof of Theorem~\ref{thm:Qcoh(P(E))}. The point is that if $\sF \in \Qcoh(\P(\sE))$ is perfect, then so is each $\sG_m \in \Qcoh(\P(\sE))$, since $q^*$ and $q_*$ preserve perfect complexes (the latter because $q$ is smooth and proper).

\subsection{Projective bundle formula}

From Theorem~\ref{thm:Perf(P(E))} and Lemma~\ref{lem:additive} we deduce:

\begin{coro}\label{cor:E(P(E))}
Let $X$ be a derived Artin stack, $\sE$ a locally free $\sO_X$-module of rank $n+1$, $n\ge 0$, and $q : \P(\sE) \to X$ the associated projective bundle. Then for any additive invariant $E$, there is a canonical isomorphism
\[
E(\P(\sE)) \simeq \bigoplus_{k=0}^{n} E(X)
\]
induced by the functors $q^*(-)\otimes \sO(-k) : \Perf(X) \to \Perf(\P(\sE))$.
\end{coro}

\section{The blow-up formula}\label{sec:blowup}

\subsection{Virtual Cartier divisors}

Recall from~\cite{KhanBlowup} that a \emph{virtual (effective) Cartier divisor} on a derived Artin stack $X$ is a quasi-smooth closed immersion $i : D \to X$ of virtual codimension~$1$. For any such $i : D \to X$, there is a canonical exact triangle
\[
\sO_X(-D) \to \sO_X \to i_*(\sO_D),
\]
where $\sO_X(-D)$ is a locally free sheaf of rank $1$, equipped with a canonical isomorphism $i^*(\sO_X(-D)) \simeq \sN_{D/X}$ (see~\cite[3.2.3 and 3.2.9]{KhanBlowup}).

\begin{lemm}\label{lem:i^*i_*(O_D)}
Let $X$ be a derived Artin stack and $i : D \to X$ a virtual Cartier divisor. Then there is a canonical isomorphism
\[
i^*i_*(\sO_D) \simeq \sO_D \oplus \sN_{D/X}[1].
\]
\end{lemm}

\begin{proof}
Applying $i^*$ to the exact triangle above (and rotating), we get the exact triangle
\[
\sO_D \to i^*i_*(\sO_D) \to \sN_{D/X}[1].
\]
The map $\sO_D \to i^*i_*(\sO_D)$ is induced by the natural transformation $i^*(\eta) : i^* \to i^*i_*i^*$ (where $\eta$ is the adjunction unit), so by the triangle identities it has a retraction given by the co-unit map $i^*i_*(\sO_D) \to \sO_D$. In other words, the triangle splits.
\end{proof}

\subsection{Grothendieck duality}

Let $i : Z \to X$ be a quasi-smooth closed immersion of derived Artin stacks. The functor $i_*$ admits a right adjoint $i^!$, which for formal reasons can be computed by the formula
\[
i^!(-) \simeq i^*(-) \otimes \omega_{D/X},
\]
where $\omega_{D/X} := i^!(\sO_X)$ is called the \emph{relative dualizing sheaf}. See~\cite[Corollary 6.4.2.7]{SAG-20180204}. When $i$ is a virtual Cartier divisor, $\omega_{D/X}$ can be computed as follows:

\begin{prop}[Grothendieck duality]\label{prop:omega_D/X}
Let $X$ be a derived Artin stack. Then for any virtual Cartier divisor $i : D \to X$, there is a canonical isomorphism
\[
\sN_{D/X}^\vee[-1] \isoto \omega_{D/X}
\]
of perfect complexes on $D$. In particular, there is a canonical identification\linebreak$i^! \simeq i^*(-) \otimes \sN_{D/X}^\vee[-1]$.
\end{prop}

\begin{proof}
Write $\sL := \sO_X(-D)$ and consider again the exact triangle $\sL \to \sO_X \to i_*(\sO_D)$. By the projection formula, this can be refined to an exact triangle of natural transformations $\id\otimes\sL \to \id \to i_*i^*$, or, passing to right adjoints, an exact triangle $i_*i^! \to \id \to \id \otimes \sL^\vee$. In particular we get the exact triangle
\begin{equation}\label{eq:divisor sequence dual}
i_*i^!(\sO_X) \to \sO_X \to \sL^\vee.
\end{equation}
The associated map $\sL^\vee[-1] \to i_*i^!(\sO_X)$ gives by adjunction a canonical morphism
\[
\sN_{D/X}^\vee[-1] \simeq i^*(\sL^\vee)[-1] \to i^!(\sO_X),
\]
which we claim is invertible. By fpqc descent and the fact that $i^!$ commutes with the operation $f^*$, for any morphism $f$~\cite[Proposition~6.4.2.1]{SAG-20180204}, we may assume that $X$ is affine. In this case the functor $i_*$ is conservative, so it will suffice to show that the canonical map
\[
i_*\left(\sN_{D/X}^\vee\right)[-1] \to i_*i^!(\sO_X)
\]
is invertible. Considering again the triangle $\sF \otimes \sL \to \sF \to i_*i^*(\sF)$ above and taking $\sF = \sL^\vee$, we get the exact triangle
\begin{equation*}
\sO_X \to \sL^\vee \to i_*i^*(\sL^\vee) \simeq i_*\left(\sN_{D/X}^\vee\right),
\end{equation*}
since $\sL$ is invertible. Comparing with~\eqref{eq:divisor sequence dual} yields the claim.
\end{proof}

\subsection{Semi-orthogonal decomposition on \texorpdfstring{$\Qcoh(\Bl_{Z/X})$}{Qcoh(Bl)}}

In this subsection we prove:

\begin{theo}\label{thm:Qcoh(Bl)}
Let $X$ be a derived Artin stack and $i : Z \to X$ a quasi-smooth closed immersion of virtual codimension $n\ge 1$. Let $\widetilde{X} = \Bl_{Z/X}$ and consider the quasi-smooth blow-up square~\eqref{eq:qsbl square}
\[
\begin{tikzcd}
D \ar{r}{i_D}\ar{d}{q} & \widetilde{X} \ar{d}{p}
\\
Z \ar{r}{i} & X
\end{tikzcd}
\]
Then we have:
\begin{enumerate}\romanenumi
\item \label{item:Qcoh(Bl)/fully faithful}
The functor $p^* : \Qcoh(X) \to \Qcoh(\widetilde{X})$ is fully faithful. We denote its essential image by $\bD(0) \subset \Qcoh(\widetilde{X})$.

\item \label{item:Qcoh(Bl)/fully faithful 2}
The functor $(i_D)_*(q^*(-) \otimes \sO(-k)) : \Qcoh(Z) \to \Qcoh(\widetilde{X})$ is fully faithful, for each $1\le k\le n-1$. We denote its essential image by $\bD(-k) \subset \Qcoh(\widetilde{X})$.

\item \label{item:Qcoh(Bl)/orthogonal}
For each $1\le k\le n-1$, the full stable subcategory $\bD(-k) \subset \Qcoh(\widetilde{X})$ is right orthogonal to each of $\bD(0),\,\ldots,\,\bD(-k+1)$.

\item \label{item:Qcoh(Bl)/generates}
The stable $\infty$-category $\Qcoh(\widetilde{X})$ is generated by the full subcategories $\bD(0)$, $\bD(-1)$, \ldots, $\bD(-n+1)$ under finite colimits and finite limits. In particular, the sequence $(\bD(0), \bD(-1),\,\ldots,\,\bD(-n+1))$ forms a semi-orthogonal decomposition of $\Qcoh(\widetilde{X})$.
\end{enumerate}
\end{theo}

\subsubsection{Proof of~\eqref{item:Qcoh(Bl)/fully faithful}}

The claim is that for any $\sF \in \Qcoh(X)$, the unit map $\sF \to p_*p^*(\sF)$ is invertible. By fpqc descent we may reduce to the case where $X$ is affine and $i$ fits in a cartesian square of the form~\eqref{eq:quasi-smooth}. Since $\Qcoh(X)$ is then generated under colimits and finite limits by $\sO_X$, and $p_*$ commutes with colimits since $p$ is quasi-compact and schematic (\ref{sssec:Qcoh functoriality}), we may assume that $\sF = \sO_X$. In other words, it suffices to show that the canonical map $\sO_X \to p_*(\sO_{\widetilde{X}})$ is invertible.
\[
\begin{tikzcd}
D \ar{r}{i_D}\ar{d}{q} & \widetilde{X}\ar{d}{p}
\\
Z \ar{r}{i} & X
\end{tikzcd}
\qquad
\begin{tikzcd}
\P^{n-1} \ar{r}\ar{d} & \Bl_{\{0\}/\A^n}\ar{d}{p_0}
\\
\{0\} \ar{r}{i_0} & \A^n,
\end{tikzcd}
\]
Since the left-hand square is the (derived) base change of the right-hand square along the morphism $f : X \to \A^n$, it follows that the map $\sO_X \to p_*(\sO_{\widetilde{X}})$ is the inverse image of the canonical map $\sO_{\A^n} \to (p_0)_*(\sO_{\Bl_{\{0\}/\A^n}})$. Thus we reduce to the case where $i$ is the immersion $\{0\} \to \A^n$. This is well-known, see~\cite[Example.~VII]{SGA6}.

\subsubsection{Proof of~\eqref{item:Qcoh(Bl)/fully faithful 2}}

It suffices to show the unit map $\sF \to q_*(i_D)^!(i_D)_*q^*(\sF)$ is invertible for all $\sF \in \Qcoh(Z)$. As in the previous claim we may assume $X$ is affine and that $\sF = \sO_Z$. Using Proposition~\ref{prop:omega_D/X}, the canonical identification $\sN_{D/\widetilde{X}} \simeq \sO_D(1)$, and Lemma~\ref{lem:i^*i_*(O_D)}, the unit map is identified with
\[
\sO_Z
\to q_*\big((i_D)^*(i_D)_*(\sO_D) \otimes \sO_D(-1)\big)[-1]
\]
Since $q : D \to Z$ is the projection of the projective bundle $\P(\sN_{Z/X})$, it follows from Proposition~\ref{prop:Serre} that we have identifications $q_*(\sO_D(-1)) \simeq 0$ and $q_*(\sO_D) \simeq \sO_Z$, under which the map in question is the identity.

\subsubsection{Proof of~\eqref{item:Qcoh(Bl)/orthogonal}}

To see that $\bD(-k)$ is right orthogonal to $\bD(0)$, observe that by Theorem~\ref{thm:Qcoh(P(E))}, the mapping space
\[
\Maps(p^*(\sF_X), (i_D)_*(q^*(\sF_Z) \otimes \sO(-k)))
\simeq \Maps(q^*i^*(\sF_X), q^*(\sF_Z) \otimes \sO(-k))
\]
is contractible for every $\sF_X \in \Qcoh(X)$ and $\sF_Z \in \Qcoh(Z)$.

To see that $\bD(-k)$ is right orthogonal to $\bD(-k')$, for $1\le k'< k$, consider the mapping space
\[
\Maps\big((i_D)_*(q^*(\sF_Z) \otimes \sO(-k')), (i_D)_*(q^*(\sF'_Z) \otimes \sO(-k))\big),
\]
for $\sF_Z,\sF'_Z \in \Qcoh(Z)$. Using fpqc descent and base change for $(i_D)_*$ against $f^*$ for any morphism $f : U \to \widetilde{X}$, we may reduce to the case where $X$ is affine. Since $\Qcoh(Z)$ is then generated under colimits and finite limits by $\sO_Z$, we may assume that $\sF_Z = \sF'_Z = \sO_Z$. Then we have
\begin{align*}
\Maps\big((i_D)_*(\sO(-k')), (i_D)_*(\sO(-k))\big) &\simeq \Maps\big((i_D)^*(i_D)_*(\sO(-k')), \sO(-k)\big)\\
&\simeq \Maps\big(\sO(-k') \oplus \sO(-k'+1)[1], \sO(-k)\big)
\end{align*}
by Lemma~\ref{lem:i^*i_*(O_D)} and the projection formula, and this space is contractible by Theorem~\ref{thm:Qcoh(P(E))}.

\subsubsection{Proof of~\eqref{item:Qcoh(Bl)/generates}}

Denote by $\bD$ the full subcategory of $\Qcoh(\widetilde{X})$ generated by $\bD(0)$, $\bD(-1),\,\ldots,\linebreak\bD(-n+1)$ under finite colimits and finite limits. The claim is that the inclusion $\bD \subseteq \Qcoh(\widetilde{X})$ is an equality. Note that $\sO_{\widetilde{X}} \in \bD(0) \subset \bD$ and $(i_D)_*(\sO_D(-k)) \in \bD(-k) \subset \bD$ for $1\le k\le n-1$. Consider the exact triangle $\sO_{\widetilde{X}}(-D) \to \sO_{\widetilde{X}} \to (i_D)_*(\sO_D)$ and recall that $\sO_{\widetilde{X}}(-D) \simeq \sO_{\widetilde{X}}(1)$. Tensoring with $\sO(-k)$ and using the projection formula, we get the exact triangle
\[
\sO_{\widetilde{X}}(-k+1) \to \sO_{\widetilde{X}}(-k) \to (i_D)_*(\sO_D(-k))
\]
for each $1\le k\le n-1$. Taking $k=1$ we deduce that $\sO_{\widetilde{X}}(-1) \in \bD$. Continuing recursively we find that $\sO_{\widetilde{X}}(-k) \in \bD$ for all $1\le k\le n-1$.

Now let $\sF \in \Qcoh(\widetilde{X})$. Denote by $\sG_0 \in \Qcoh(\widetilde{X})$ the cofibre of the co-unit map $p^*p_*(\sF) \to \sF$. Note that $\sG_0$ is right orthogonal to $\bD(0)$. For $1\le m\le n-1$ define $\sG_m$ recursively by the exact triangles
\[
(i_D)_*\big(q^*q_*((i_D)^!(\sG_{m-1}) \otimes \sO(m)) \otimes \sO(-m)\big)
\xrightarrow{\mrm{counit}} \sG_{m-1} \to \sG_m.
\]
Just as in the proof of Theorem~\ref{thm:Qcoh(P(E))}, a simple induction argument shows that each $\sG_m$ is right orthogonal to all of the subcategories $\bD(0),\,\ldots,\,\bD(m-1)$. We now claim that $\sG_{n-1}$ is zero; it will follow by recursion that $\sF$ belongs to $\bD$, as desired.

Since the objects $\sG_k$ are stable under base change, we may use fpqc descent and base change to assume that $X$ is affine. Moreover we may assume that $i: Z\to X$ fits in a cartesian square of the form~\eqref{eq:quasi-smooth}. By~\cite[3.3.6]{KhanBlowup}, $p : \widetilde{X} \to X$ factors through a quasi-smooth closed immersion $i' : \widetilde{X} \to \P^{n-1}_X$. Recall from Lemma~\ref{lem:O(n+1)} that there is a canonical isomorphism $\colim_{J \subsetneq [n-1]} \sO(\abs{J}) \simeq \sO(n)$ in $\Qcoh(\P^{n-1}_X)$. Applying $(i')^*$, we get 
\[
\colim_{J \subsetneq [n-1]} \sO_{\widetilde{X}}(\abs{J}) \simeq \sO_{\widetilde{X}}(n) 
\]
in $\Qcoh(\widetilde{X})$.
In particular, every $\sO_{\widetilde{X}}(k)$ belongs to $\bD$ for all $k\in\bZ$. Recall also that we may find a map $\bigoplus_\alpha \sO(d_\alpha)[n_\alpha] \to i'_*(\sG_{n-1})$ which is surjective on all homotopy groups (Lemma~\ref{lem:surjection from O(m)'s}). By adjunction this corresponds to a map $\bigoplus_\alpha \sO(d_\alpha)[n_\alpha] \to \sG_{n-1}$ (which is also surjective on homotopy groups). But the source belongs to $\bD$, and the target is right orthogonal to $\bD$, so this map is null-homotopic. Thus $\sG_{n-1}$ is zero.

\subsection{Proof of Theorem~\ref{thm:Perf(Bl)}}\label{ssec:proof of Perf(Bl)}

We now deduce Theorem~\ref{thm:Perf(Bl)} from Theorem~\ref{thm:Qcoh(Bl)}. First note that the fully faithful functor $\sF \mapsto p^*(\sF)$ of Theorem~\ref{thm:Qcoh(Bl)}
\eqref{item:Qcoh(P(E))/fully faithful} preserves perfect complexes and therefore restricts to a fully faithful functor $\Perf(X) \to \Perf(\Bl_{Z/X})$. This shows Theorem~\ref{thm:Perf(Bl)}$\MK$\eqref{item:Perf(Bl)/fully faithful 1}.

Similarly, part
\eqref{item:Perf(Bl)/fully faithful 2} follows from the fact that the functors $q^*$ and $(i_D)_*$ preserve perfect complexes. For the latter, this is because $i_D$ is quasi-smooth (and hence of finite presentation and of finite tor-amplitude).

For part
\eqref{item:Perf(Bl)/SOD} we argue again as in the proof of Theorem~\ref{thm:Qcoh(Bl)}$\MK$\eqref{item:Qcoh(Bl)/generates}. The point is that if $\sF \in \Qcoh(\Bl_{Z/X})$ is perfect, then so is each $\sG_m \in \Qcoh(\P(\sE))$, since $q^*$, $q_*$, $(i_D)_*$ and $(i_D)^!$ all preserve perfect complexes. For the latter this follows from Proposition~\ref{prop:omega_D/X}.
\goodbreak


\subsection{Blow-up formula}\label{ssec:blowup/additive}


By Theorem~\ref{thm:Perf(Bl)} and Lemma~\ref{lem:additive} we get:

\begin{coro}\label{cor:E(Bl_{Z/X})}
Let $X$ be a derived Artin stack and $i : Z \to X$ a quasi-smooth closed immersion of virtual codimension $n\ge 1$. Then for any additive invariant $E$, there is a canonical isomorphism
\[
E\left(\Bl_{Z/X}\right) \simeq E(X) \oplus \bigoplus_{k=1}^{n-1} E(Z).
\]
\end{coro}

\subsubsection{Proof of Theorem~\ref{thm:blow-up descent}}\label{sssec:proof of blow-up descent}

Combine Corollaries~\ref{cor:E(Bl_{Z/X})} and~\ref{cor:E(P(E))} $\left(\text{with }\sE = \sN_{Z/X}\right)$.

\section{The cdh topology}\label{sec:cdh}

\subsection{Definitions}

The following notion was introduced by Voevodsky~\cite{VoevodskyCdh} for noetherian schemes:

\begin{defi}\label{defn:cdh square}
%\leavevmode
Suppose given a cartesian square $Q$ of algebraic spaces
\begin{equation}\label{eq:cdh square}
\begin{tikzcd}
B \ar{r}\ar{d} & Y \ar{d}{p}
\\
A \ar{r}{e} & X.
\end{tikzcd}
\end{equation}
\begin{enumerate}\romanenumi
\item \label{item:cdh square/Nis square}
We say that $Q$ is a \emph{Nisnevich square} if $e$ is an open immersion, and $p$ is an étale morphism inducing an isomorphism $(Y\setminus B)_\red \simeq (X\setminus A)_\red$.
\item \label{item:cdh square/proper cdh square}
We say that $Q$ is a \emph{proper cdh square}, or \emph{abstract blow-up square}, if $e$ is a closed immersion of finite presentation, and $p$ is a proper morphism inducing an isomorphism $(Y\setminus B)_\red \simeq (X\setminus A)_\red$.
\item \label{item:cdh square/cdh square}
We say that $Q$ is a \emph{cdh square} if it is either a Nisnevich square or a proper cdh square.
\end{enumerate}
\end{defi}


Given any class of commutative squares of algebraic spaces, we say that a presheaf satisfies \emph{descent} for this class if it sends all such squares to homotopy cartesian squares, and the empty scheme to a terminal object. In case of the three classes considered in
Definition~\ref{defn:cdh square}, it follows from a theorem of Voevodsky~\cite[Corollary~5.10]{VoevodskyCD} that descent in this sense is equivalent to Čech descent with respect to the associated Grothendieck topology.

\begin{exam}\label{exam:Nis descent for E localizing}
Every localizing invariant $E$ satisfies Nisnevich descent when regarded as a presheaf on quasi-compact quasi-separated algebraic spaces with $E(X)\linebreak= E(\Perf(X))$. This is essentially due to Thomason~\cite{ThomasonTrobaugh} and in the asserted generality is a consequence of the study of compact generation properties of the $\infty$-categories\;$\Qcoh(X)$ carried out by Bondal--Van den Bergh~\cite{BondalVDB}.
\end{exam}

\begin{exam}\label{exam:qsbl square cl}
Any quasi-smooth blow-up square~\eqref{eq:qsbl square} induces a proper cdh square
\[
\begin{tikzcd}
\P\left(\sN_{Z/X}|_{Z_\cl}\right) \ar{r}\ar{d} & \left(\Bl_{Z/X}\right)_\cl \ar{d}
\\
Z_\cl \ar{r} & X_\cl
\end{tikzcd}
\]
on underlying classical algebraic spaces.
\end{exam}

\begin{exam}\label{exam:closed square}
Consider the class of proper cdh squares~\eqref{eq:cdh square} where the proper morphism $p$ is a closed immersion (with quasi-compact open complement). The associated Grothendieck topology is the same as the one generated by \emph{closed squares}, i.e. cartesian squares as in~\eqref{eq:cdh square} such that $e$ and $p$ are closed immersions, $e$ is of finite presentation and $p$ has quasi-compact open complement, and $A \sqcup Y \to X$ is surjective on underlying topological spaces.
\end{exam}

\begin{exam}\label{exam:nil-immersion}
Note that for any algebraic space $X$, the square
\[
\begin{tikzcd}
\initialobj \ar{r}\ar{d} & X_\red \ar{d}
\\
\initialobj \ar{r} & X
\end{tikzcd}
\]
is a closed square as in Example~\ref{exam:closed square}.
\end{exam}

\subsection{A cdh descent criterion}

\begin{theo}\label{thm:cdh criterion}
Let $\sF$ be a presheaf on the category $\bC$ of algebraic spaces, with values in a stable $\infty$-category. Then $\sF$ satisfies cdh descent if and only if it satisfies the following conditions:
\begin{enumerate}\romanenumi
\item\label{item:cdh criterion/reduced}
It sends the empty scheme to a zero object.
\item\label{item:cdh criterion/Nis}
It sends Nisnevich squares to cartesian squares.
\item\label{item:cdh criterion/closed}
It sends closed squares to cartesian squares.
\item\label{item:cdh criterion/qsbl}
For every $X\in\bC$ and every quasi-smooth closed immersion $Z \to X$, it sends the square (Example~\ref{exam:qsbl square cl})
\[
\begin{tikzcd}
\P\left(\sN_{Z/X}|_{Z_\cl}\right) \ar{r}\ar{d} & \left(\Bl_{Z/X}\right)_\cl \ar{d}
\\
Z_\cl \ar{r} & X
\end{tikzcd}
\]
to a cartesian square.
\end{enumerate} Moreover, the same holds if $\bC$ is replaced by the full subcategory of
%%les alphenumi sont introduites en inlinelist  à rechercher !
\begin{enumerate}\alphenumi
\item quasi-compact quasi-separated (qcqs) algebraic spaces,
\item schemes,
\item or qcqs schemes.
\end{enumerate}
\end{theo}

\begin{rema}\label{rem:trivial extension to dass}
Any presheaf $\sF$ on algebraic spaces can be trivially extended to derived algebraic spaces, by setting $\Gamma(X, \sF) = \Gamma(X_\cl, \sF)$ for every derived algebraic space $X$. The condition~
\eqref{item:cdh criterion/qsbl} in Theorem~\ref{thm:cdh criterion} is equivalent to requiring this extension to satisfy descent for quasi-smooth blow-up squares~\eqref{eq:qsbl square}.
\end{rema}

\begin{exam}\label{exam:cdh=localizing+nil+closed}
Let $E$ be a localizing invariant of stable $\infty$-categories. Then it satisfies Nisnevich descent on qcqs algebraic spaces (Example~\ref{exam:Nis descent for E localizing}) and quasi-smooth blow-up descent (Theorem~\ref{thm:blow-up descent}). Assume that $E$ also satisfies \emph{derived nilpotent invariance}, i.e., that the canonical map $E(X) \to E(X_\cl)$ is invertible for every derived algebraic space $X$. Then the condition~\eqref{item:cdh criterion/qsbl} in Theorem~\ref{thm:cdh criterion} holds. Therefore, $E$ satisfies cdh descent if and only if it satisfies closed descent. Moreover, by Nisnevich descent it suffices to consider closed squares of affine schemes.
\end{exam}

\begin{exam}\label{exam:A1 closed}
In the presence of $\A^1$-homotopy invariance, the Morel--Voevodsky localization theorem~\cite[Theorem~3.2.21]{MorelVoevodsky} provides the following sufficient condition for closed descent. Let $\sF$ be an $\A^1$-invariant Nisnevich sheaf on the category of algebraic spaces. Suppose that, for every algebraic space $S$, its restriction $\sF_S$ to the site of \emph{smooth} algebraic spaces over $S$ is stable under arbitrary base change. That is, for every morphism of algebraic spaces $f : T \to S$, the canonical map $f^*(\sF_S) \to \sF_T$ is invertible. Then $\sF$ satisfies closed descent. This follows immediately from the closed base change formula (cf.~\cite[Proposition~3.3.2]{KhanLocalization}).
\end{exam}

\begin{rema}
Let $E$ be a localizing invariant and suppose that it is moreover \emph{truncating} in the sense of~\cite{LandTamme}. That is, if $R$ is a connective $\sE_1$-ring spectrum and $\Mod_R^\perf$ denotes the stable $\infty$-category of left $R$-modules, then the canonical map $E(\Mod^\perf_R) \to E(\Mod^\perf_{\pi_0(R)})$ is invertible. Then Land--Tamme have recently proven that $E$ has closed descent, at least if we restrict to \emph{noetherian} algebraic spaces (see Step~1 in the proof of~\cite[Theorem.~A.2]{LandTamme}).
\end{rema}

\begin{rema}\label{rem:variants of cdh criterion}
There are a few variants of Theorem~\ref{thm:cdh criterion} with the same proof. For example:
\begin{enumerate}\romanenumi
\item On the category of (qcqs) schemes, descent with respect to the \emph{rh} topology (generated by Zariski squares and proper cdh squares) can be checked with the same criteria, except that Nisnevich squares are replaced by Zariski squares in condition
\eqref{item:cdh criterion/Nis}.
\item If we do not assume either Nisnevich or Zariski descent, descent for the proper cdh topology is still equivalent to conditions~\eqref{item:cdh criterion/reduced},
\eqref{item:cdh criterion/closed}, and~\eqref{item:cdh criterion/qsbl}, as long as we restrict to a full subcategory of algebraic spaces or schemes which satisfy Thomason's \emph{resolution property}. For example, this holds on the category of quasi-projective schemes.
\item \label{item:variants of cdh criterion/stacks}
One can extend the criterion to qcqs Artin stacks as follows. The definition of Nisnevich square extends without modification (cf.~\cite[Subsection~2.3]{HoyoisKrishna}). In the definition of proper cdh square, we add the requirement that the proper morphism $p$ is representable (cf. op. cit.). Then the criterion of Theorem~\ref{thm:cdh criterion} holds for stacks \emph{which admit the resolution property Nisnevich-locally}, see (Subsubsection~\ref{sssec:cdh stacks}). This condition is relatively mild. For example, many quotient stacks have the resolution property (\cite[Lemma~2.4]{ThomasonEquivariant}, \cite[Example~7.5]{HallRydhPerfect}). By the Nisnevich-local structure theorem of Alper--Hall--Rydh~\cite[Theorem~2.9]{HoyoisKrishna}, any stack with linearly reductive and almost multiplicative stabilizers satisfies the resolution property Nisnevich-locally.
\end{enumerate}
\end{rema}

\subsection{Proof of Theorem~\ref{thm:cdh criterion}}

Since Nisnevich squares, closed squares, and quasi-smooth blow-up squares are all cdh squares, the conditions are clearly necessary. Conversely suppose that $\sF$ is a presheaf satisfying the conditions and let $Q$ be a proper cdh square of algebraic spaces of the form
\begin{equation}\label{eq:proper cdh in proof}
\begin{tikzcd}
E \ar{r}\ar{d} & Y \ar{d}{p}
\\
Z \ar{r}{i} & X.
\end{tikzcd}
\end{equation}
It will suffice to show that the induced square $\Gamma(Q, \sF)$ is homotopy cartesian.

\subsubsection{}\label{sssec:cdh criterion proof/blow-up}

Assume first that $Q$ is a blow-up square, i.e., that $Y = \Bl_{Z/X}$ is the blow-up of $X$ centred in $Z$ (and $E = \P(\sC_{Z/X})$ is the projectivized normal cone). By Nisnevich descent we may assume that $X$ satisfies the resolution property (e.g. $X$ is affine). Since $i : Z \to X$ is of finite presentation, the ideal of definition $\sI \subset \sO_X$ is of finite type. Thus by the resolution property there exists a surjection $u : \sE \to \sI$ with $\sE$ a locally free $\sO_X$-module of finite rank. Denote by $V = \bV_X(\sE) = \Spec_X(\Sym_{\sO_X}(\sE))$ the associated vector bundle and $0 : X \to V$ the zero section. The $\sO_X$-module homomorphism $u : \sE \to \sI \subset \sO_X$ induces a section of $V$, whose derived zero-locus $\widetilde{Z}$ fits in the homotopy cartesian square
\[
\begin{tikzcd}
\widetilde{Z} \ar{r}{\widetilde{i}}\ar{d} & X \ar{d}{u}
\\
X \ar{r}{0} & V.
\end{tikzcd}
\]
By construction, $\widetilde{i} : \widetilde{Z} \to X$ is a quasi-smooth closed immersion and there is a canonical morphism $Z \to \widetilde{Z}$ which induces an isomorphism $Z \simeq \widetilde{Z}_\cl$. Regarding $\sF$ as a presheaf on derived algebraic spaces as in Remark~\ref{rem:trivial extension to dass}, the square $\Gamma(Q,\sF)$ now factors as follows:
\[
\begin{tikzcd}
\Gamma(X, \sF) \ar{r}\ar{d} & \Gamma(\widetilde{Z}, \sF) \ar{d}\ar[equals]{r} & \Gamma(Z, \sF) \ar{ldd}
\\
\Gamma\left(\Bl_{\widetilde{Z}/X}, \sF\right) \ar{r}\ar{d} & \Gamma\left(\P(\sN_{\widetilde{Z}/X}\right), \sF) \ar{d}
\\
\Gamma\left(\Bl_{Z/X}, \sF\right) \ar{r} & \Gamma\left(\P\left(\sC_{Z/X}\right), \sF\right)
\end{tikzcd}
\]
The upper square is induced by a quasi-smooth blow-up square, hence is cartesian. The lower square is induced by a closed square, hence is also cartesian. Therefore it follows that the outer composite square is also cartesian. This shows that $\sF$ satisfies descent for blow-up squares.

\subsubsection{}\label{sssec:cdh criterion proof/blow-up (X-Z)-admissible}

Slightly more generally, suppose that $Y = \Bl_{Z'/X}$ is a blow-up centred in some closed immersion $Z' \to X$ with $\abs{Z'} \subseteq \abs{Z}$ on underlying topological spaces, and let $E' \to Y$ denote the exceptional divisor. Since $\sF$ is invariant under nilpotent extensions (Example~\ref{exam:nil-immersion}) we may assume that $i' : Z' \to X$ actually factors through a closed immersion $Z' \to Z$ (see Example~\ref{exam:nil-immersion}). Applying descent for blow-up squares (Subsubsection~\ref{sssec:cdh criterion proof/blow-up}), it will suffice to show that $\sF$ satisfies descent for the square
\[
\begin{tikzcd}
E' \ar{r}\ar{d} & E \ar{d}
\\
Z' \ar{r} & Z.
\end{tikzcd}
\]
Note that the blow-up $\Bl_{Z'/Z}$ is equipped with a canonical closed immersion into $E$ so that $E' \to E$ and $\Bl_{Z'/Z} \to Z$ form a closed covering. Applying closed descent and descent for the blow-up square associated to $Z' \to Z$ (Subsubsection~\ref{sssec:cdh criterion proof/blow-up}), we~conclude.

\subsubsection{}

For the case of an arbitrary proper cdh square, we recall the standard argument using Raynaud--Gruson's technique of \emph{platification par éclatements}~\cite[I, Corollary~5.7.12]{RaynaudGruson} to reduce to the case considered above (see e.g.~\cite[Subsection~5.2]{KerzStrunkTamme}).

\goodbreak
\begingroup
\begin{constr}\label{constr}
Let $Q$ be a proper cdh square of the form~\eqref{eq:proper cdh in proof}. Assume that $X$ is quasi-compact and quasi-separated and that the open subspace $X\setminus Z$ is quasi-compact and dense in $X$. Then there exists a proper cdh square $Q'$ sitting above $Q$ such that the composite $Q' \circ Q$
\[
\begin{tikzcd}
E' \ar{r}\ar{d}\ar[phantom]{rd}{\scriptstyle Q'} & Y' \ar{d}
\\
E \ar{r}\ar{d}\ar[phantom]{rd}{\scriptstyle Q} & Y \ar{d}{p}
\\
Z \ar[swap]{r}{i} & X
\end{tikzcd}
\]
is of the form considered in (Subsubsection~\ref{sssec:cdh criterion proof/blow-up (X-Z)-admissible}). That is, $Y'$ is a blow-up of $X$ centred in some closed immersion $Z' \to X$ with $\abs{Z'} \subseteq \abs{Z}$. This follows from Raynaud--Gruson just as in the proof of~\cite[Corollary~2.4]{HoyoisKrishna}.
\end{constr}
\endgroup

Let $Q$ be a proper cdh square of the form~\eqref{eq:proper cdh in proof}. Since $\sF$ is a Nisnevich sheaf, we may assume that $X$ is quasi-compact and quasi-separated. Since $i : Z \to X$ is of finite presentation, its open complement $U = X\setminus Z$ is quasi-compact. Using closed descent, we can ensure that $U$ is dense in $X$. Now apply the construction above to get a proper cdh square $Q'$ such that $Q'\circ Q$ is of the form considered in (Subsubsection~\ref{sssec:cdh criterion proof/blow-up (X-Z)-admissible}). Applying the construction again, this time to $Q'$, we end up with a third square $Q''$ such that the composite $Q'' \circ Q'$ is also of the form considered in (Subsubsection~\ref{sssec:cdh criterion proof/blow-up (X-Z)-admissible}). Then we know that $\Gamma(Q'\circ Q, \sF)$ and $\Gamma(Q''\circ Q', \sF)$ are both homotopy cartesian. It follows that the square $\Gamma(Q', \sF)$ is also homotopy cartesian (since $\sF$ takes values in a stable $\infty$-category, it suffices to check that the induced map on homotopy fibres is invertible), and hence so is $\Gamma(Q, \sF)$.

\subsubsection{}\label{sssec:cdh stacks}

We now discuss the extension to stacks mentioned in Remark~\ref{rem:variants of cdh criterion}$\MK$\eqref{item:variants of cdh criterion/stacks}. The precise statement is as follows. Let $\bC$ be a category of qcqs Artin stacks such that
\begin{enumerate}\alphenumi
\item every stack $X\in\bC$ admits a Nisnevich atlas by stacks with the resolution property;
\item for every stack $X\in\bC$ and every blow-up $Y \to X$, the qcqs Artin stack $Y$ also belongs to $\bC$.
\end{enumerate}
Then the statement of Theorem~\ref{thm:cdh criterion} holds for presheaves on $\bC$.

The proof for the case of a blow-up square (Subsubsection~\ref{sssec:cdh criterion proof/blow-up}) has been presented in such a way that it holds \emph{mutatis mutandis} under the above assumptions. The argument of~\cite[Claim~5.3]{KerzStrunkTamme} also goes through, using descent for closed squares and blow-up squares, to deal with the slightly more general case where $Y = \Bl_{Z'/X}$ is a blow-up centred in some closed immersion $Z' \to X$ that factors through $Z$. To reduce a general proper cdh square to that case, we use Rydh's extension of Raynaud--Gruson~\cite[Theorem~2.2]{HoyoisKrishna}. First, closed descent allows us to assume that $X\setminus Z$ is dense in $X$. Then we apply Rydh--Raynaud--Gruson just as in the proof of~\cite[Corollary~2.4]{HoyoisKrishna}. The only difference with the case of schemes or algebraic spaces is that in general we get a \emph{sequence} of $(X\setminus Z)$-admissible blow-ups $\tilde{X} \to X$ which factors through $p : Y \to X$. The addition of a simple induction is then the only modification required to run the same argument.

\subsection{Homotopy invariant K-theory}\label{ssec:kh/kh}

\subsubsection{}

For any qcqs algebraic space $X$, its homotopy invariant K-theory spectrum is given by the formula
\begin{equation}\label{eq:Gamma(X, Lhtp(K))}
\Gamma(X, \KH) = \colim_{[n] \in \bDelta^{\op}} \K(X \times \A^n).
\end{equation}
In other words, $\Gamma(X, \KH)$ is the geometric realization of the simplicial diagram $\K(X \times \A^\bullet)$, where $\A^\bullet$ is regarded as a cosimplicial scheme in the usual way (see e.g.~\cite[p.~45]{MorelVoevodsky}). This extends the usual definition~\cite{WeibelKH,ThomasonTrobaugh}, and is a way to formally impose the property of $\A^1$-homotopy invariance: for any qcqs algebraic space $X$, the projection $p : X \times \A^1 \to X$ induces an isomorphism of spectra
\[
p^* : \Gamma(X, \KH) \to \Gamma(X \times \A^1, \KH).
\]

\subsubsection{}\label{sssec:KH derived}

As the previous paragraph makes sense when $X$ is derived, we may regard $\KH$ as a presheaf on qcqs derived algebraic spaces. Given a Nisnevich square of the form~\eqref{eq:cdh square}, Nisnevich descent for K-theory (Example~\ref{exam:Nis descent for E localizing}) yields homotopy cartesian squares of spectra
\[
\begin{tikzcd}
\K(X \times \A^n) \ar{r}\ar{d} & \K(A \times \A^n) \ar{d}
\\
\K(Y \times \A^n) \ar{r} & \K(B \times \A^n)
\end{tikzcd}
\]
for every $[n] \in \bDelta^{\op}$. Passing to the colimit over $n$, we deduce that $\KH$ also satisfies Nisnevich descent. We have:

\begin{theo}\label{thm:KH nil}
For every qcqs derived algebraic space $S$, the canonical morphism of spectra
\[
\Gamma(S, \KH) \to \Gamma(S_\cl, \KH)
\]
is invertible.
\end{theo}

\begin{proof}
By Nisnevich descent, we may as well assume $S$ is an affine derived scheme. Let $\KH_S$ denote the restriction of $\KH$ to the site of affine derived schemes that are smooth and of finite presentation over $S$. This is still an $\A^1$-homotopy invariant Nisnevich sheaf, and it is equipped with a canonical morphism
\[
\K^{\cn}_S \to \K_S \to \KH_S,
\]
where $\K^{\cn}_S$ and $\K_S$ are the respective restrictions of connective and non-connective K-theory. By Cisinski, this morphism exhibits $\KH_S$ as the \emph{Bott periodization} of the $\A^1$-localization of $\K^{\cn}_S$, i.e., the periodization with respect to the Bott element $b \in \K_1(\bG_{m,S})$ (the proof is the same as in the case where $S$ is classical~\cite[Corollary~2.12]{CisinskiKH}). It follows from this description that for any morphism of affine derived schemes $f : T \to S$, there is a canonical isomorphism $f^*(\KH_S) \simeq \KH_T$, where $f^*$ denotes the functor of inverse image of $\A^1$-invariant Nisnevich sheaves. Indeed, we reduce to checking the same property for $\K^{\cn}_S$, which is clear as this is identified up to Zariski localization with the group completion of the presheaf $\coprod_n \mrm{BGL}_{n,S}$.

In particular, we get a canonical isomorphism $i^*(\KH_S) \simeq \KH_{S_\cl}$, where $i : S_\cl \to S$ is the inclusion of the underlying classical scheme. Moreover, $i^*$ induces an equivalence between the $\infty$-categories of $\A^1$-invariant Nisnevich sheaves on $S$ and $S_\cl$, respectively, by~\cite[Corollary~1.3.5]{KhanThesis}. We deduce that the canonical morphism
\[
\Gamma(S, \KH)
\simeq \Gamma(S, \KH_S)
\to \Gamma(S_\cl, \KH_{S_\cl})
\simeq \Gamma(S_\cl, \KH)
\]
is invertible.
\end{proof}

\subsubsection{Proof of Theorem~\ref{thm:KH cdh descent}}\label{sssec:proof of KH cdh descent}

We use the criterion of Theorem~\ref{thm:cdh criterion}. Condition~
\eqref{item:cdh criterion/reduced} is obvious. Nisnevich descent (condition~
\eqref{item:cdh criterion/Nis}) was verified above (Subsubsection~\ref{sssec:KH derived}). For condition~
\eqref{item:cdh criterion/qsbl}, it will suffice by Theorem~\ref{thm:KH nil} and Remark~\ref{rem:trivial extension to dass} to show that $\KH$ sends quasi-smooth blow-up squares of derived algebraic spaces to homotopy cartesian squares. This follows from the same property for K-theory (Theorem~\ref{thm:blow-up descent}) using the formula~\eqref{eq:Gamma(X, Lhtp(K))} (just as in the proof of Nisnevich descent). For closed descent (condition~
\eqref{item:cdh criterion/closed}), we may restrict our attention to closed squares of \emph{affine} schemes (by Nisnevich descent). This is classical, see~\cite[Exer.~9.11$\MK$(f)]{ThomasonTrobaugh} or~\cite[Corollary~4.10]{WeibelKH}. Alternatively, it follows from the criterion of Example~\ref{exam:A1 closed}.

\begin{rema}
By continuity for $\KH$ (e.g.~\cite[Theorem.~4.9$\MK$(5)]{HoyoisKrishna}), once we have descent for proper cdh squares as in Definition~\ref{defn:cdh square}$\MK$\eqref{item:cdh square/proper cdh square}, we can immediately drop the finite presentation hypothesis on $e$.
\end{rema}


\begin{thebibliography}{MMMMMM}
\bibitem[BGI71]{SGA6} Berthelot, Pierre; Grothendieck, Alexandre; Illusie, Luc Séminaire de Géométrie Algébrique du Bois Marie – 1966-67 – Théorie des intersections et théorème de Riemann–Roch – (SGA 6), Lecture Notes in Mathematics, 225, Springer, 1971 \href{https://ahl.centre-mersenne.org/item/AHL_2020__3__1091_0/#SGA6}{Publisher reference}.
\bibitem[BGT13]{BlumbergGepnerTabuada} Blumberg, Andrew J.; Gepner, David; Tabuada, Gonçalo A universal characterization of higher algebraic K-theory, Geom. Topol., Volume 17 (2013) no. 2, pp. 733-838 \href{https://ahl.centre-mersenne.org/item/AHL_2020__3__1091_0/#BlumbergGepnerTabuada}{Publisher reference}.
\bibitem[BK89]{BondalKapranov} Bondal, Aleksei Igorevich; Kapranov, Mikhail Mikhailovich Representable functors, Serre functors, and mutations, Izv. Ross. Akad. Nauk SSSR, Ser. Mat., Volume 53 (1989) no. 6, pp. 1183-1205 \href{https://ahl.centre-mersenne.org/item/AHL_2020__3__1091_0/#BondalKapranov}{Publisher reference}.
\bibitem[BM12]{BlumbergMandellTHH} Blumberg, Andrew J.; Mandell, Michael A. Localization theorems in topological Hochschild homology and topological cyclic homology, Geom. Topol., Volume 16 (2012) no. 2, pp. 1053-1120 \href{https://ahl.centre-mersenne.org/item/AHL_2020__3__1091_0/#BlumbergMandellTHH}{Publisher reference}.
\bibitem[BS20]{BerghSchnuerer} Bergh, Daniel; Schnürer, Olaf M Conservative descent for semi-orthogonal decompositions, Adv. Math., Volume 360 (2020), 106882, 39 pages \href{https://ahl.centre-mersenne.org/item/AHL_2020__3__1091_0/#BerghSchnuerer}{Publisher reference}.
\bibitem[BVdB03]{BondalVDB} Bondal, Aleksei Igorevich; Van den Bergh, Michel Generators and representability of functors in commutative and noncommutative geometry, Mosc. Math. J., Volume 3 (2003) no. 1, pp. 1-36 \href{https://ahl.centre-mersenne.org/item/AHL_2020__3__1091_0/#BondalVDB}{Publisher reference}.
\bibitem[Cis13]{CisinskiKH} Cisinski, Denis-Charles Descente par éclatements en K-théorie invariante par homotopie, Ann. Math., Volume 177 (2013) no. 2, pp. 425-448 \href{https://ahl.centre-mersenne.org/item/AHL_2020__3__1091_0/#CisinskiKH}{Publisher reference}.
\bibitem[GR17]{GaitsgoryRozenblyum} Gaitsgory, Dennis; Rozenblyum, Nick A Study in Derived Algebraic Geometry: Volumes I and II, Mathematical Surveys and Monographs, 221.1-221.2, American Mathematical Society, 2017 \href{https://ahl.centre-mersenne.org/item/AHL_2020__3__1091_0/#GaitsgoryRozenblyum}{Publisher reference}.
\bibitem[Hae04]{Haesemeyer} Haesemeyer, Christian Descent properties of homotopy K-theory, Duke Math. J., Volume 125 (2004) no. 3, pp. 589-619 \href{https://ahl.centre-mersenne.org/item/AHL_2020__3__1091_0/#Haesemeyer}{Publisher reference}.
\bibitem[HK19]{HoyoisKrishna} Hoyois, Marc; Krishna, Amalendu Vanishing theorems for the negative K-theory of stacks, Ann. K-Theory, Volume 4 (2019) no. 3, pp. 439-472 \href{https://ahl.centre-mersenne.org/item/AHL_2020__3__1091_0/#HoyoisKrishna}{Publisher reference}.
\bibitem[HLP14]{HalpernLeistnerPreygel} Halpern-Leistner, Daniel; Preygel, Anatoly Mapping stacks and categorical notions of properness (2014) (\url{https://arxiv.org/abs/1402.3204}) \href{https://ahl.centre-mersenne.org/item/AHL_2020__3__1091_0/#HalpernLeistnerPreygel}{Publisher reference}.
\bibitem[Hoy16]{HoyoisCdh} Hoyois, Marc Cdh descent in equivariant homotopy K-theory (2016) (\url{https://arxiv.org/abs/1604.06410}) \href{https://ahl.centre-mersenne.org/item/AHL_2020__3__1091_0/#HoyoisCdh}{Publisher reference}.
\bibitem[HR17]{HallRydhPerfect} Hall, Jack; Rydh, David Perfect complexes on algebraic stacks, Compos. Math., Volume 153 (2017) no. 11, pp. 2318-2367 \href{https://ahl.centre-mersenne.org/item/AHL_2020__3__1091_0/#HallRydhPerfect}{Publisher reference}.
\bibitem[Kha16]{KhanThesis} Khan, Adeel A. Motivic homotopy theory in derived algebraic geometry, Ph. D. Thesis, Universität Duisburg-Essen, Germany (2016) (Available at \url{https://www.preschema.com/thesis/}) \href{https://ahl.centre-mersenne.org/item/AHL_2020__3__1091_0/#KhanThesis}{Publisher reference}.
\bibitem[Kha19]{KhanLocalization} Khan, Adeel A. The Morel–Voevodsky localization theorem in spectral algebraic geometry, Geom. Topol., Volume 23 (2019) no. 7, pp. 3647-3685 \href{https://ahl.centre-mersenne.org/item/AHL_2020__3__1091_0/#KhanLocalization}{Publisher reference}.
\bibitem[KR18a]{KhanBlowup} Khan, Adeel A.; Rydh, David Virtual Cartier divisors and blow-ups (2018) (\url{https://arxiv.org/abs/1802.05702}) \href{https://ahl.centre-mersenne.org/item/AHL_2020__3__1091_0/#KhanBlowup}{Publisher reference}.
\bibitem[KR18b]{KrishnaRavi} Krishna, Amalendu; Ravi, Charanya Algebraic K-theory of quotient stacks, Ann. K-Theory, Volume 3 (2018) no. 2, pp. 207-233 \href{https://ahl.centre-mersenne.org/item/AHL_2020__3__1091_0/#KrishnaRavi}{Publisher reference}.
\bibitem[KST18]{KerzStrunkTamme} Kerz, Moritz; Strunk, Florian; Tamme, Georg Algebraic K-theory and descent for blow-ups, Invent. Math., Volume 211 (2018) no. 2, pp. 523-577 \href{https://ahl.centre-mersenne.org/item/AHL_2020__3__1091_0/#KerzStrunkTamme}{Publisher reference}.
\bibitem[LT18]{LandTamme} Land, Markus; Tamme, Georg On the K-theory of pullbacks (2018) (\url{https://arxiv.org/abs/1808.05559}) \href{https://ahl.centre-mersenne.org/item/AHL_2020__3__1091_0/#LandTamme}{Publisher reference}.
\bibitem[Lur09]{HTT} Lurie, Jacob Higher topos theory, Annals of Mathematics Studies, Princeton University Press, 2009 no. 170 \href{https://ahl.centre-mersenne.org/item/AHL_2020__3__1091_0/#HTT}{Publisher reference}.
\bibitem[Lur12]{HA-20170918} Lurie, Jacob Higher algebra, 2012 (Preprint, version of 2017-09-18 available at \url{https://www.math.ias.edu/~lurie/papers/HA.pdf}) \href{https://ahl.centre-mersenne.org/item/AHL_2020__3__1091_0/#HA-20170918}{Publisher reference}.
\bibitem[Lur16]{SAG-20180204} Lurie, Jacob Spectral algebraic geometry (2016) (Preprint, version of 2018-02-03, available at \url{https://www.math.ias.edu/~lurie/papers/SAG-rootfile.pdf}) \href{https://ahl.centre-mersenne.org/item/AHL_2020__3__1091_0/#SAG-20180204}{Publisher reference}.
\bibitem[MV99]{MorelVoevodsky} Morel, Fabien; Voevodsky, Vladimir A 1 -homotopy theory of schemes, Publ. Math., Inst. Hautes Étud. Sci., Volume 90 (1999) no. 1, pp. 45-143 \href{https://ahl.centre-mersenne.org/item/AHL_2020__3__1091_0/#MorelVoevodsky}{Publisher reference}.
\bibitem[Orl92]{OrlovSOD} Orlov, Dmitri Olegovich Projective bundles, monoidal transformations, and derived categories of coherent sheaves, Izv. Ross. Akad. Nauk, Ser. Mat., Volume 56 (1992) no. 4, pp. 852-862 \href{https://ahl.centre-mersenne.org/item/AHL_2020__3__1091_0/#OrlovSOD}{Publisher reference}.
\bibitem[RG71]{RaynaudGruson} Raynaud, Michel; Gruson, Laurent Critères de platitude et de projectivité: Techniques de « platification »  d’un module, Invent. Math., Volume 13 (1971) no. 1-2, pp. 1-89 \href{https://ahl.centre-mersenne.org/item/AHL_2020__3__1091_0/#RaynaudGruson}{Publisher reference}.
\bibitem[Tho87]{ThomasonEquivariant} Thomason, Robert Wayne Equivariant resolution, linearization, and Hilbert’s fourteenth problem over arbitrary base schemes, Adv. Math., Volume 65 (1987) no. 1, pp. 16-34 \href{https://ahl.centre-mersenne.org/item/AHL_2020__3__1091_0/#ThomasonEquivariant}{Publisher reference}.
\bibitem[Tho93a]{ThomasonProjectiveBundle} Thomason, Robert Wayne Les K-groupes d’un fibré projectif, Algebraic K-theory and algebraic topology (Lake Louise. Proceedings of the NATO Advanced Study Institute, Lake Louise, Alberta, Canada, December 12–16, 1991 (NATO ASI Series. Series C. Mathematical and Physical Sciences), Volume 407, Kluwer Academic Publishers, 1993, pp. 243-248 \href{https://ahl.centre-mersenne.org/item/AHL_2020__3__1091_0/#ThomasonProjectiveBundle}{Publisher reference}.
\bibitem[Tho93b]{ThomasonBlowup} Thomason, Robert Wayne Les K-groupes d’un schéma éclaté et une formule d’intersection excédentaire, Invent. Math., Volume 112 (1993) no. 1, pp. 195-215 \href{https://ahl.centre-mersenne.org/item/AHL_2020__3__1091_0/#ThomasonBlowup}{Publisher reference}.
\bibitem[TT90]{ThomasonTrobaugh} Thomason, Robert Wayne; Trobaugh, Thomas Higher algebraic K-theory of schemes and of derived categories, The Grothendieck Festschrift. Vol. III (Progress in Mathematics), Volume 88, Springer, 1990, pp. 247-435 \href{https://ahl.centre-mersenne.org/item/AHL_2020__3__1091_0/#ThomasonTrobaugh}{Publisher reference}.
\bibitem[TV08]{HAG2} Toën, Bertrand; Vezzosi, Gabriele Homotopical algebraic geometry II: Geometric stacks and applications, 902, American Mathematical Society, 2008 \href{https://ahl.centre-mersenne.org/item/AHL_2020__3__1091_0/#HAG2}{Publisher reference}.
\bibitem[Voe10a]{VoevodskyCD} Voevodsky, Vladimir Homotopy theory of simplicial sheaves in completely decomposable topologies, J. Pure Appl. Algebra, Volume 214 (2010) no. 8, pp. 1384-1398 \href{https://ahl.centre-mersenne.org/item/AHL_2020__3__1091_0/#VoevodskyCD}{Publisher reference}.
\bibitem[Voe10b]{VoevodskyCdh} Voevodsky, Vladimir Unstable motivic homotopy categories in Nisnevich and cdh-topologies, J. Pure Appl. Algebra, Volume 214 (2010) no. 8, pp. 1399-1406 \href{https://ahl.centre-mersenne.org/item/AHL_2020__3__1091_0/#VoevodskyCdh}{Publisher reference}.
\bibitem[Wei89]{WeibelKH} Weibel, Charles A. Homotopy algebraic K-theory, Algebraic K-theory and algebraic number theory (Honolulu, HI, 1987) (Contemporary Mathematics), Volume 83, American Mathematical Society, 1989, pp. 461-488 \href{https://ahl.centre-mersenne.org/item/AHL_2020__3__1091_0/#WeibelKH}{Publisher reference}.
\end{thebibliography}
\end{document}
